\documentclass[11pt,a4paper]{article}
% Estrutura e apresentação adaptadas do template fornecido pelo autor.
\usepackage[utf8]{inputenc}
\usepackage[T1]{fontenc}
\usepackage[brazil,english]{babel}
\usepackage{lmodern}
\usepackage[margin=1in]{geometry}
\usepackage{authblk}
\usepackage{amsmath,amssymb,amsfonts}
\usepackage{graphicx,booktabs,xcolor,hyperref,cite,microtype,fancyhdr}
\usepackage{longtable,array,tabularx,enumitem}
\usepackage{xurl}
\hypersetup{colorlinks=true,linkcolor=blue!50!black,citecolor=blue!50!black,urlcolor=blue!50!black,pdftitle={Do riso impresso ao kkkk},pdfauthor={Murillo Edson de Carvalho Souza}}
\pagestyle{fancy}\fancyhf{}
\fancyhead[L]{\small Do riso impresso ao \textit{kkkk}}
\fancyhead[R]{\small\thepage}
\setlength{\headheight}{15pt}
\renewcommand{\headrulewidth}{0.4pt}
\setlength{\emergencystretch}{3em}
\newcolumntype{Y}{>{\raggedright\arraybackslash}X}
\newcommand{\forma}[1]{\textit{#1}}
\newcommand{\repo}{\url{https://github.com/murilloed/historia-riso-escrito-portugues}}
\author{\textbf{Murillo Edson de Carvalho Souza}}
\date{}
\title{\textbf{Do riso impresso ao \textit{kkkk}: mudança linguística, iconicidade e evidência documental na história das representações do riso em português}\\[0.6em]\large Auditoria de alegações etimológicas e proposta de amostragem histórica reproduzível}
\newcommand{\CorpusSize}{121}
\newcommand{\AnnualAdded}{96}
\newcommand{\AnnualCovered}{110}
\newcommand{\AnnualMissing}{0}

\begin{document}
\selectlanguage{brazil}\maketitle
\begin{abstract}
Investiga-se em que condições representações históricas do riso em português podem ser relacionadas à convenção digital \forma{kkkk}. O estudo articula linguística histórica, iconicidade e crítica documental em coleta exploratória não probabilística. Os 16 endereços indicados foram organizados em 13 conjuntos de recursos de acesso; bibliotecas, periódicos, repositórios e documentos complementares foram rastreados por URL e identificador. Os inventários incluem documentos desde 1645, com 40 identificadores de fonte anteriores a 1900, sem cobertura anual contínua desse período. O núcleo contém \CorpusSize{} registros: 96 rubricas descritivas e 25 representações miméticas ou marcadores digitais, com cobertura anual de 1900 a 2009. A seleção por acervo e década compreende 135 vínculos e 129 fontes por identificador; 65 vínculos têm confirmação no acervo indicado e 70 em fontes alternativas. Dois exemplares específicos foram conferidos sem ocorrência localizada, com alternativas literárias de período registradas separadamente. As 66 indicações iniciais receberam parecer; nenhuma das 22 transcrições alegadas foi autenticada integralmente. Confirmam-se testemunhos de \forma{quiá}, \forma{cá}, \forma{ah} e \forma{quá}; uma dissertação de 2008 documenta \forma{kkkk} em pesquisa situada em 2007.2. Os dados mostram diversidade de soluções, mas não estabelecem origem absoluta ou cadeia de transmissão entre o riso oitocentista e a escrita digital. Os cálculos de amostragem e detecção constituem planejamento prospectivo, sem estimação populacional para a coleta realizada. Conclui-se que semelhança não demonstra filiação e lacuna documental não demonstra independência. Dados, localizadores e scripts públicos permitem auditar contagens, procedência e limites das interpretações.
\end{abstract}
\noindent\textbf{Palavras-chave:} riso escrito; onomatopeia; português brasileiro; linguística histórica; comunicação digital; amostragem documental.
\begin{otherlanguage}{english}
\begin{abstract}
This study examines whether historical written laughter in Portuguese can be linked to digital \textit{kkkk}. An exploratory, non-probability documentary design combines historical linguistics, iconicity and source criticism. Sixteen requested entry points are organized into thirteen operational access-resource groups. The inventories include documents dating back to 1645 and 40 source identifiers predating 1900, without continuous annual coverage of that period. The core contains \CorpusSize{} records: 96 descriptive laughter annotations and 25 mimetic representations or digital markers, covering every year from 1900 to 2009. Archive-by-decade selection comprises 135 selection links and 129 source identifiers: 65 links were confirmed in the indicated archive and 70 through alternative sources. Two fixed newspaper items yielded no attestation in visual inspection; literary alternatives for their periods are recorded separately. All 66 initial bibliographic leads were audited, but none of the 22 alleged quotations was fully authenticated. Historical forms and digital \textit{kkkk} are documented, including a 2008 dissertation reporting research conducted in 2007. These findings establish neither an absolute origin nor a transmission chain from nineteenth-century writing. Sampling and detection calculations are prospective scenarios, not population estimates from this collection. Resemblance alone does not demonstrate descent, and documentary gaps do not demonstrate independence. Public data and scripts support review of evidence, counts and limitations.
\end{abstract}
\end{otherlanguage}
\noindent\textbf{Keywords:} written laughter; onomatopoeia; Brazilian Portuguese; historical linguistics; digital communication; documentary sampling.
\clearpage
\section{Introdução}
Uma gargalhada pode aparecer na escrita como \forma{ah! ah! ah!}, \forma{quiá-quiá}, \forma{cá-cá}, \forma{quá-quá}, \forma{hahaha}, \forma{rsrsrs} ou \forma{kkkk}. Na comunicação digital, o repertório também inclui emoticons, emojis, sequências heterogêneas de letras, imagens e respostas verbais como \forma{rindo}. A convergência funcional dessas formas é evidente: todas podem participar da representação ou da negociação do humor. Sua relação histórica exige outra espécie de evidência. Compartilhar um referente ou uma motivação sonora não demonstra que uma forma descenda de outra.

A pergunta central deste artigo é: \emph{quais evidências permitem distinguir continuidade histórica, convergência icônica e inovação gráfica na representação do riso em português?} A questão ganha relevância quando explicações de circulação pública ligam a escrita digital brasileira a passagens de José de Alencar, Machado de Assis ou Eça de Queirós, sem mostrar os elos intermediários. Um registro oitocentista pode ser um antecedente documental da representação escrita do riso, mas isso não o converte automaticamente em origem etimológica de uma convenção digital.

O problema inverso é igualmente importante. Seria incorreto rejeitar essa genealogia afirmando que onomatopeias, por definição, não possuem história, não podem ser transmitidas ou nunca se integram à gramática. O caráter imitativo de uma expressão não a impede de ser aprendido, convencionalizado, emprestado, lexicalizado ou transformado. A literatura recente mostra a participação da experiência linguística na percepção de iconicidade \cite{iida2024}. A alternativa entre ``evolução real'' e ``mera onomatopeia'' precisa ser substituída por dimensões que podem coexistir.

O material que motivou o estudo reúne explicações conversacionais sobre mudança fonética, semântica e gramatical; exemplos como \forma{olho}, \forma{você} e \forma{vilão}; comparações entre idiomas; alegações sobre a internet brasileira dos anos 1990 e 2000; e panoramas decenais de livros, jornais e quadrinhos. Inclui frases apresentadas como citações literais, às vezes acompanhadas de datas precisas, mas sem edição, página ou reprodução documental. Essa combinação torna necessária uma crítica de fontes: a precisão aparente de uma data não substitui a recuperação do documento.

O objetivo geral é construir uma interpretação linguística e documental das representações do riso, preservando as perguntas do material de partida e examinando criticamente suas respostas. Os objetivos específicos são: definir mudança, derivação, parentesco, iconicidade e coincidência; revisar exemplos etimológicos; recuperar testemunhos com localizadores; auditar cronologias; fundamentar estatisticamente a ampliação do corpus; e disponibilizar instrumentos de reprodução. A repetição de uma afirmação em várias respostas não produz novas fontes independentes.

Emprega-se \emph{português brasileiro} para a variedade investigada, sem pressupor que uma única gramática normativa explique todos os usos digitais. ``Gramática da língua brasileira'', no material inicial, é interpretada como interesse pela organização do português em uso no Brasil. Distinguem-se \emph{sorriso}, expressão facial que pode ser silenciosa, \emph{riso}, categoria mais ampla, e \emph{gargalhada}, manifestação vocal intensa. O objeto é o riso escrito: um marcador digital de humor não comprova que o escrevente tenha rido fisicamente.

Os fundamentos de linguística e de fonética e fonologia apoiam-se em obras de referência \cite{martelotta2008,seara2015}. Estudos de 2024 e 2025 complementam os fundamentos em iconicidade e comunicação digital \cite{iida2024,badrakh2024,kempe2025}.

\section{Problema historiográfico e hipóteses concorrentes}
\subsection{Anterioridade, descendência e equivalência funcional}
Há quatro proposições diferentes: uma forma aparece antes de outra; duas formas representam riso; usuários de uma tiveram contato com a outra; uma convenção foi transformada a partir da outra. As duas primeiras são compatíveis com múltiplas histórias. A terceira exige evidência de circulação; a quarta, uma explicação de transmissão apoiada em documentação e contexto.

O enunciado ``o \forma{kkkk} nasceu em \textit{Til}, em 1872'' funde essas proposições. A presença de \forma{quiá} não equivale à presença do caractere repetido \forma{k}. Por outro lado, ``o \forma{quiá} morreu nos livros'' também excede a evidência sem pesquisa suficientemente ampla de cartas, periódicos, teatro, quadrinhos e práticas escolares. Obras impressas podem funcionar como canais de transmissão; uma tradição escrita não depende de passagem obrigatória pela fala cotidiana.

Consideram-se três hipóteses: H1, transmissão ou adaptação de convenções preexistentes; H2, recriação independente motivada por iconicidade e propriedades gráficas; H3, origem múltipla ou convergente, com transmissão local, reinvenção e difusão em redes. Nenhuma recebe exclusividade por definição. A semelhança entre \forma{quiá}, \forma{cá} e o nome da letra \forma{k} torna relações plausíveis, mas plausibilidade não demonstra causalidade.

\subsection{Primeira ocorrência encontrada e origem absoluta}
Uma ocorrência datada estabelece que a forma estava disponível naquele documento ou na interação nele reproduzida. Oferece, se o testemunho for confiável, um limite superior para a emergência: a inovação ocorreu até aquela data. Não informa necessariamente onde, por quem ou quanto antes. ``Primeira ocorrência encontrada neste levantamento'' difere de ``primeira ocorrência da língua''.

A conservação de registros é desigual. Conversas privadas de ICQ ou MSN, canais de IRC, BBS e salas públicas nem sempre foram arquivados. Um log preservado exige validação de data, integridade e procedência. A ausência de \forma{kkkk} em um conjunto de conversas entre 1995 e 1997, caso recuperado, seria resultado sobre aquele conjunto, não prova de inexistência no Brasil. Sem corpus identificado, a própria alegação de ausência permanece sem base verificável.

\section{Fundamentos linguísticos e revisão dos conceitos}
\subsection{Mudança diacrônica não se reduz à derivação}
Mudança linguística designa transformações em sistemas e práticas de linguagem ao longo do tempo. Pode afetar pronúncia, contrastes, estruturas morfológicas, construções sintáticas, significados, usos pragmáticos e convenções de escrita. ``Evolução'' significa aqui mudança histórica, sem supor progresso, aperfeiçoamento ou seleção biológica literal. A diversidade dos mecanismos é central nos estudos históricos \cite{campbell2021,martelotta2008}.

Na morfologia sincrônica, \emph{raiz} e \emph{radical} dizem respeito à estrutura de unidades analisadas em determinado sistema, com definições dependentes da teoria. A \emph{etimologia} investiga a origem histórica. \emph{Derivação morfológica} é um mecanismo de formação de palavras, como a combinação de base e afixo; não é requisito para que uma palavra tenha história. A descendência de um item pode envolver erosão e reanálise que tornem a estrutura ancestral opaca. A capacidade de receber sufixos não é teste necessário nem suficiente de continuidade.

Correspondências sonoras regulares ajudam a reconstruir relações; analogia, empréstimo e mudança semântica também precisam ser considerados. Não é necessário observar diretamente cada geração intermediária. A reconstrução se apoia na convergência de evidências, não numa gravação contínua de todos os estágios. A metáfora do ``DNA dos vocábulos'' tem limites: palavras não possuem material genético biológico e podem atravessar fronteiras entre línguas e suportes.

\subsection{Fonética, fonologia, grafia e nome de letra}
A fonética examina propriedades articulatórias, acústicas e perceptivas; a fonologia investiga a organização dos contrastes sonoros. A grafia tem convenções próprias e não é transcrição acústica transparente. O sinal \forma{k} pode ativar o nome da letra, normalmente \forma{cá} em português, mas não é tecnicamente correto chamar \forma{cá} de ``fonema da letra'': é uma sílaba composta por segmentos. Letra, som e fonema precisam ser distinguidos \cite{seara2015}.

A consoante /k/ é uma oclusiva velar surda, articulada com fechamento pelo dorso da língua na região do véu palatino. A vogal acompanhante não integra esse fonema. As letras \forma{c}, diante de \forma{a}, e o dígrafo \forma{qu}, em muitos contextos, podem representar /k/. Disso não resulta identidade fonética integral entre \forma{cá}, \forma{quá} e \forma{quiá}. Nem toda gargalhada contém oclusivas velares. A classificação de /b/ como dental ou alveolar, presente no material inicial, é incorreta: /b/ é bilabial.

Síncope é perda no interior de uma forma; apócope, perda no final; aférese, perda no início. A supressão inicial em \forma{você} $>$ \forma{cê} é aférese, não apócope da primeira sílaba. Metátese, assimilação, redução e palatalização são processos distintos. A ``lei do menor esforço'' pode ser uma intuição sobre tendências articulatórias, mas não explica sozinha a propagação social de inovações.

\subsection{Onomatopeia, iconicidade e convencionalização}
Onomatopeia é representação linguisticamente organizada de som. Seu caráter icônico é uma relação percebida de semelhança entre forma e conteúdo. A semelhança é seletiva e culturalmente mediada: sistemas fonológicos, repertórios gráficos e práticas de leitura participam da interpretação. Iconicidade e convenção não são incompatíveis. Iida e Akita \cite{iida2024} mostram que familiaridade e experiência na língua influenciam julgamentos de iconicidade.

Uma expressão imitativa pode ser criada num episódio e depois repetida por outros usuários. Pode ser emprestada, adaptada ou incorporada a processos morfológicos. Portanto, ``é onomatopeia'' não implica ``não tem etimologia''. Ocorrências interjetivas podem apresentar autonomia prosódica ou sintática relativa, o que não significa ausência de organização linguística. Uma convenção digital pode ser histórica mesmo quando sua transformação é principalmente gráfica.

\emph{Coincidência fonética} é uma hipótese sobre a relação entre formas: elas se parecem sem que a semelhança decorra de descendência ou empréstimo demonstrável. Pode haver coincidência casual ou convergência favorecida pelo mesmo estímulo acústico. Não se devem fundir coincidência e onomatopeia numa categoria. Uma palavra não onomatopaica pode coincidir com outra; duas onomatopeias podem ter relação histórica.

\begin{table}[htbp]\centering\small
\caption{Dimensões que substituem a oposição rígida entre ``evolução real'' e ``onomatopeia''.}
\begin{tabularx}{\textwidth}{p{0.19\textwidth}YY}\toprule
Dimensão & Questão pertinente & Inferência indevida\\\midrule
História & Há transmissão ou empréstimo documentável? & Semelhança prova descendência.\\
Motivação & Há semelhança percebida com um som? & Iconicidade exclui convenção.\\
Morfologia & A forma participa de construções ou famílias? & Sem afixação não há história.\\
Pragmática & Que ação realiza na interação? & Todo marcador implica riso físico.\\
Suporte & Que recursos e restrições são relevantes? & O teclado determina a inovação.\\
Documentação & Qual é a fonte, a data e o localizador? & Não encontrar prova ausência.\\\bottomrule
\end{tabularx}\end{table}

\subsection{Continuidade lexical em três idiomas}
O português \forma{olho} exemplifica continuidade lexical a partir do latim \forma{oculum}, convencionalmente citado no acusativo. Sua história envolve redução e transformação consonantal, com palatalização no percurso românico. Não se reduz à ``sonorização de oc-''. \forma{Voar}, relacionado a \forma{volare}, e \forma{lua}, relacionado a \forma{luna}, ilustram perdas intervocálicas em contextos pertinentes, sem uma receita irrestrita para todas as palavras \cite{ilari2018,campbell2021}.

No inglês, \forma{good} remonta ao inglês antigo \forma{g\=od}. O exemplo combina continuidade e mudança de pronúncia, sem grande alteração gráfica. Sua trajetória não deve ser atribuída mecanicamente à Grande Mudança Vocálica: histórias particulares incluem encurtamentos e condicionamentos. No francês, \forma{eau} é relacionado ao latim \forma{aqua}, com forte divergência gráfica e sonora; a etimologia é registrada pela Académie française \cite{gooddict,academieeau}. Nenhum exemplo exige reconhecimento consciente da origem pelos falantes atuais.

\begin{table}[htbp]\centering\small
\caption{Três exemplos de continuidade lexical.}
\begin{tabularx}{\textwidth}{p{0.14\textwidth}p{0.24\textwidth}Y}\toprule
Idioma & Relação histórica & Cautela\\\midrule
Português & \forma{oculum} $>$ \forma{olho} & Fundamenta-se na história românica, não só no mesmo referente.\\
Inglês & \forma{g\=od} $>$ \forma{good} & Mudanças particulares não se reduzem a uma única lei geral.\\
Francês & \forma{aqua} $>$ \forma{eau} & A relação pode persistir com pouca semelhança superficial.\\\bottomrule
\end{tabularx}\end{table}

\subsection{\forma{Você}: gramaticalização e história social}
\forma{Vossa mercê} $>$ \forma{você} é redução e reorganização histórica de uma expressão de tratamento, associada à pronominalização. A expressão-fonte contém possessivo e nome; a forma resultante desempenha funções de referência ao interlocutor. A análise envolve categoria, distribuição e significado social. Não é um exemplo direto de palavra latina que apenas preservou seu radical: o ponto de partida imediato é uma construção portuguesa.

A cronologia ``vossa mercê no XVI, vosmecê no XIX, você no XX e cê no XXI'' é inadequada. Monte \cite{monte2019} documenta \forma{você} em 1638, refutando seu suposto nascimento no século XX. As variantes não se sucedem como degraus exclusivos: coexistem, distribuem-se regionalmente e adquirem especializações pragmáticas. Eventos como a chegada da corte em 1808 podem integrar hipóteses sociais, mas não substituem documentos datados.

A concordância de \forma{você vai} conserva propriedades da construção nominal de tratamento em sua integração ao sistema pronominal. A referência é ao interlocutor, enquanto a concordância padrão tem forma de terceira pessoa. Não é uma ``anomalia''. A distribuição de \forma{tu} e \forma{você} também varia no Brasil. Sua história articula estrutura linguística, polidez, circulação social e organização colonial \cite{monte2019}.

\subsection{\forma{Vilão}: semântica, estrutura e acepções coexistentes}
\forma{Vilão}, historicamente relacionado a \forma{villanus} e à base \forma{villa}, é pertinente à pejoração: denominações sociais podem adquirir avaliações negativas. Entretanto, uma narrativa única de ``camponês livre'', ``pessoa rude'' e ``antagonista'' não descreve todas as épocas e regiões. Estatutos sociais e acepções dependem do contexto; sentidos antigos e recentes podem coexistir.

A explicação formal \forma{-us} $>$ \forma{-ão} é incorreta. A história das terminações e da nasalidade precisa considerar a base em \forma{-anu-} e a evolução românica. No plano gramatical, \forma{vilão} pode funcionar como substantivo e adjetivo: um sentido moral mais abstrato não obriga a mudança de classe. Categoria, significado e forma não mudam necessariamente em bloco \cite{ilari2018,campbell2021}.

\subsection{Comparações internacionais e falsas etimologias}
O latido em \forma{au-au}, \forma{woof} ou \forma{bow-wow} e o choro em \forma{buá} ou \forma{ouin} mostram diversidade de convenções, mas não demonstram sozinhos independência histórica. Famílias linguísticas diferentes não impedem empréstimos, inclusive de formas icônicas. É preciso investigar cada item.

Um contraexemplo ao material inicial é o japonês \forma{chikutaku}: o verbete reunido em Kotobank identifica relação com o inglês \forma{ticktack} \cite{kotobank}. Apresentá-lo como prova de criação isolada entre Japão e Europa é impróprio. Sem análise especializada, não se adotam como fatos uma origem imitativa do mandarim \forma{gou}, um suposto inglês antigo \forma{gaw}, uma leitura egípcia \forma{iu-iu} ou uma origem onomatopaica consensual do latim \forma{basium}. Etimologia remota incerta não prova imitação independente.

Formas como o tailandês \forma{555} e a repetição coreana de uma consoante ilustram recursos gráficos diferentes: nome de numeral e sinal consonantal. A comparação tipológica não exige afirmar isolamento cultural. A transliteração ``chau-chau'' atribuída à forma coreana no material inicial não é empregada como descrição fonética.

As menções a Ælfric, Chaucer, Shakespeare e Aristófanes permanecem pistas bibliográficas, não testemunhos autenticados neste corpus. O coaxar em \textit{As rãs} não deve ser reclassificado como riso; a grafia grega contemporânea não pode ser projetada sobre a Grécia antiga sem edição e análise. A alegação do ``primeiro haha do Ocidente'' exige investigação especializada e não é adotada como resultado.

\section{Metodologia documental e organização do repositório}
Na revisão documental de 8 de novembro de 2026, cada uma das 66 indicações recebidas ganhou identificador estável e parecer individual, publicados no repositório do autor \cite{souzaRepositorio2026}. A matriz distingue autenticação de representação do riso, referência lexical ou temática, identificação bibliográfica, recuperação de exemplar específico e casos sem confirmação. As buscas literais das 22 frases e suas limitações constam em arquivos próprios. A existência de um catálogo não foi tratada como substituto de passagem primária; a pesquisa de OCR não foi descrita como leitura visual integral. O método estatístico a seguir permanece prospectivo: as novas aquisições não transformam a coleta dirigida em amostra probabilística.
\subsection{Desenho executado, escopo e material de partida}
O estudo executado é exploratório, documental e crítico, com busca dirigida e coleta não probabilística. Seu objeto empírico tem dois níveis: as alegações e supostas citações presentes no material conversacional fornecido pelo autor; e os documentos recuperados para verificá-las ou contextualizá-las. Não é uma revisão sistemática exaustiva, nem uma amostra representativa de toda a escrita portuguesa entre 1872 e 2026. O recorte de 1900 a 2009 organiza a cobertura anual do núcleo histórico sugerido no material; testemunhos anteriores e posteriores funcionam como extensões documentadas.

A consulta foi realizada em 8 de novembro de 2026. Utilizaram-se buscas por autor, título, variante gráfica e ano em acervos institucionais e bibliotecas digitais, incluindo Project Gutenberg, Wikisource, o corpus de Machado de Assis mantido na Brigham Young University, a Hemeroteca Digital Brasileira e repositórios universitários. Variantes sem acento, pontuação alterada e repetições distintas foram consideradas como pistas. Resultados de buscadores serviram para descoberta; a passagem recuperada, a imagem ou o texto integral acessível fundamentaram a classificação de evidência.

O material conversacional não foi contado como ocorrência histórica. Ele é um roteiro de investigação que contém hipóteses, exemplos, contradições e transcrições ainda sem autenticação. A auditoria conserva o assunto de cada trecho, distinguindo a formulação original de sua avaliação científica. Os arquivos integrais de trabalho permanecem locais; a versão pública contém uma matriz de alegações com os elementos necessários à conferência, sem divulgar a conversa inteira ou metadados pessoais desnecessários.

\subsection{Unidades e metadados}
Distinguem-se quatro unidades. A \emph{referência bibliográfica} identifica uma obra; o \emph{documento amostral} é uma edição, número de periódico ou sessão de interação; a \emph{ocorrência} é uma passagem contextualizada; o \emph{tipo gráfico} é uma categoria, como \forma{quiá} ou \forma{kkkk}. Três passagens do mesmo romance não equivalem a três documentos independentes. Uma dissertação que reproduz um chat é o documento consultado, enquanto o chat é o testemunho reportado. A unidade não deve mudar silenciosamente entre tabelas.

Para cada testemunho registram-se identificador, autor ou periódico, título, ano da obra, ano da edição ou estudo, data da interação quando disponível, país, suporte, grafia literal, contexto, localizador, URL, data de acesso, natureza da evidência e restrições. Em PDFs, informa-se a página do arquivo e, quando conhecida, a paginação impressa. Em textos eletrônicos sem paginação estável, utiliza-se capítulo e âncora textual. Um identificador baseado em sequência exata permite reencontrar a passagem sem inventar uma página.

A grafia diplomática conserva acento, caixa e pontuação na medida permitida pelo testemunho. A normalização é um campo separado, usado para localizar variantes. ``Quá quá quá'' e ``quá-quá-quá'' podem ser agrupados numa análise de tipos, mas não são apresentados como transcrições idênticas. Um OCR que lê uma vogal incorretamente não decide a forma final: a imagem prevalece quando disponível.

\subsection{Níveis de evidência e regras de inclusão}
Adotam-se os códigos E1, fonte primária com imagem inspecionada e localizador; E2, texto primário em edição ou transcrição eletrônica consultada, sem colação integral com o impresso; E3, interação reproduzida em estudo acadêmico, com página e contexto identificados; E4, alegação secundária, paráfrase ou referência sem passagem autenticada. E1 não significa ausência de qualquer erro editorial; E2 não equivale a edição crítica. Esses códigos qualificam a cadeia de acesso, não a importância literária.

Incluem-se no inventário de atestações passagens cujo contexto sustenta leitura de riso e cujo suporte pode ser identificado. Excluem-se da contagem principal nomes comerciais, títulos de programas, siglas e sequências de letras sem função de riso demonstrada. A palavra ``risos'' numa descrição não é automaticamente rubrica de transcrição. Uma cadeia \forma{KKKK} que nomeia um empreendimento ou uma organização não é uma gargalhada digital. Tais falsos positivos são registrados quando relevantes para explicar a triagem.

As frases literais do material inicial não foram autenticadas por existirem livros com os mesmos títulos. A identificação da obra e a verificação da passagem são etapas distintas. A ausência de correspondência numa busca eletrônica também não prova falsificação: pode haver diferenças de edição, OCR ou redação. O estatuto adequado é ``não autenticada nesta coleta'', salvo quando há uma contradição documental específica.

\subsection{Rastreabilidade, downloads e disponibilidade}
O repositório do autor, \repo, concentra o manuscrito, bibliografia, dados e scripts \cite{souzaRepositorio2026}. O inventário de aquisição registra URL, resultado de acesso, tamanho e SHA-256 dos arquivos efetivamente baixados. Falhas de acesso são mantidas no histórico. Os scripts distinguem arquivos obtidos localmente de documentos consultados por serviço de leitura web; uma consulta bem-sucedida num ambiente não é apresentada como download bem-sucedido em outro.

O GitHub é infraestrutura de versionamento e disponibilização, não prova de validade da hipótese. A referência metodológica identifica o autor e o repositório; a reprodução deve usar o commit ou a versão efetivamente consultada. Links, checksums e scripts contribuem para localização e reuso segundo os princípios FAIR \cite{wilkinson2016}. Não se presume que acesso gratuito autorize republicação integral. Documentos com direitos não esclarecidos permanecem locais, com metadados públicos e instruções de obtenção. Arquivos redistribuídos conservam fonte e condições aplicáveis.

\subsection{Controle de qualidade e limites de validação}
Efetuou-se conferência contextual das passagens selecionadas, inspeção visual da capa de \textit{O Tico-Tico} e validação dos localizadores disponíveis. Uma revisão eletrônica verifica chaves bibliográficas, referências cruzadas, ausência de citações não resolvidas, contagem de registros e consistência entre tabelas geradas e arquivos de dados. A revisão é feita em uma etapa de trabalho assistida por ferramentas; não se declara dupla codificação independente ou concordância entre avaliadores que não ocorreu.

Uma ampliação deverá empregar pelo menos dois anotadores em uma subamostra, medir acordo por categoria e discutir divergências. Textos antigos exigem atenção a acentos, hifens, fontes tipográficas e ordem das colunas. A validação humana não deve limitar-se a confirmar a regex: precisa verificar a função da sequência e sua atribuição ao documento correto.

\subsection{Síntese quantitativa, inventário de bases e links de reprodução}
O quadro quantitativo reúne \CorpusSize{} testemunhos, 66 indicações auditadas e 22 citações alegadas verificadas. As 57 referências da seleção por acervo e década e os 17 títulos inéditos incluem documentos que também integram o núcleo. As tabelas indicam as unidades e os arquivos correspondentes no GitHub.

Os dezesseis endereços indicados pelo autor não equivalem a dezesseis bases independentes: o APESP aparece como portal e item individual; a BBM, como coleção específica e acervo geral; e o Orkut é conteúdo preservado na própria Wayback. O inventário conserva todos os endereços e explicita essas relações. Os recursos complementares, como bibliotecas digitais e repositórios universitários, estão identificados por URL individual no manifesto, com inventário adicional de seus domínios. Um domínio técnico, uma página de catálogo e uma obra consultada não são unidades intercambiáveis.

O manifesto histórico inclui arquivos auxiliares, metadados, código de interfaces, documentos sem ocorrência autenticada e o histórico de uma referência retirada. As quantidades de aquisições e falhas constam no quadro quantitativo. A unidade do manifesto é a entrada de aquisição.

\small\begin{longtable}{>{\raggedright\arraybackslash}p{.31\textwidth}>{\raggedright\arraybackslash}p{.13\textwidth}>{\raggedright\arraybackslash}p{.45\textwidth}}
\caption{Balanço quantitativo consolidado do levantamento. As linhas têm unidades distintas e não devem ser somadas.}\\
\toprule Indicador & Quantidade & Unidade, escopo e dados\\\midrule\endfirsthead
\toprule Indicador & Quantidade & Unidade, escopo e dados\\\midrule\endhead
Endereços indicados pelo autor & 16 & endereços. Portais, coleções, item individual e captura; há sobreposições institucionais. \href{https://github.com/murilloed/historia-riso-escrito-portugues/blob/main/data/base_registry.csv}{Dados no GitHub} \\
Núcleo de representações autenticadas & 121 & documentos/testemunhos. Inclui documentos primários e interações reportadas; exclui menções apenas lexicais. \href{https://github.com/murilloed/historia-riso-escrito-portugues/blob/main/data/attestations.csv}{Dados no GitHub} \\
Níveis do núcleo E1 / E2 / E3 & 23 / 95 / 3 & testemunhos. Imagem / texto primário eletrônico / interação reportada; partição dos 121 registros. \href{https://github.com/murilloed/historia-riso-escrito-portugues/blob/main/data/attestations.csv}{Dados no GitHub} \\
Indicações iniciais com parecer & 66 & indicações. Auditoria bibliográfica; não são 66 ocorrências positivas. \href{https://github.com/murilloed/historia-riso-escrito-portugues/blob/main/data/references_audit.csv}{Dados no GitHub} \\
Citações alegadas submetidas a verificação & 22 & alegações. Nenhuma autenticada integralmente; os resultados constam no parecer. \href{https://github.com/murilloed/historia-riso-escrito-portugues/blob/main/data/claims_verification.csv}{Dados no GitHub} \\
Referências por acervo e década & 135 & referências selecionadas. Inclui menções lexicais, rubricas e grafias miméticas, com fontes alternativas identificadas na matriz. \href{https://github.com/murilloed/historia-riso-escrito-portugues/blob/main/data/archive_decade_references.csv}{Dados no GitHub} \\
Linhas de acompanhamento por acervo/década & 137 & linhas. 135 linhas com referência e dois itens conferidos sem ocorrência; não é o universo de documentos. \href{https://github.com/murilloed/historia-riso-escrito-portugues/blob/main/data/archive_decade_matrix.csv}{Dados no GitHub} \\
Títulos inéditos selecionados & 17 & títulos. Rodada 4: 14 Brasiliana e três espaços digitais; duas fontes já incorporadas ao núcleo. \href{https://github.com/murilloed/historia-riso-escrito-portugues/blob/main/data/revision4_new_titles.csv}{Dados no GitHub} \\
Anos no inventário anual do núcleo & 155 & anos. 1872-2026; anos sem testemunho não significam ausência histórica. \href{https://github.com/murilloed/historia-riso-escrito-portugues/blob/main/data/annual_coverage.csv}{Dados no GitHub} \\
Entradas no manifesto histórico & 1635 & entradas de aquisição. Inclui apoio, metadados, recursos técnicos, falhas e referências retiradas; não documentos lidos. \href{https://github.com/murilloed/historia-riso-escrito-portugues/blob/main/data/downloads.json}{Dados no GitHub} \\
Entradas com conteúdo baixado & 1482 & entradas não vazias. Inclui HTML, PDF, imagem, JSON e código de interface; não equivale a leitura integral. \href{https://github.com/murilloed/historia-riso-escrito-portugues/blob/main/data/downloads.json}{Dados no GitHub} \\
Respostas vazias preservadas no manifesto & 2 & entradas vazias. Não contadas como documentos analisados ou resultados negativos. \href{https://github.com/murilloed/historia-riso-escrito-portugues/blob/main/data/downloads.json}{Dados no GitHub} \\
Entradas com falha de aquisição & 151 & entradas com falha. Estado registrado por identificador; não representa todas as tentativas históricas. \href{https://github.com/murilloed/historia-riso-escrito-portugues/blob/main/data/downloads.json}{Dados no GitHub} \\
Textos integrais com cópia pública & 8 & cópias públicas. Demais arquivos integrais ficam locais quando a redistribuição não foi verificada. \href{https://github.com/murilloed/historia-riso-escrito-portugues/blob/main/sources/README.md}{Dados no GitHub} \\
Referências bibliográficas & 322 & entradas BibTeX. Inclui teoria, documentos, catálogos e portais; não corresponde ao número de atestações. \href{https://github.com/murilloed/historia-riso-escrito-portugues/blob/main/artigo/references.bib}{Dados no GitHub} \\
Cobertura anual de 1900 a 2009 & 110 & anos com atestação. Intervalo de 110 anos. A busca dirigida incluiu rubricas e grafias miméticas. \href{https://github.com/murilloed/historia-riso-escrito-portugues/blob/main/data/annual_coverage.csv}{Dados no GitHub} \\
Conjuntos de recursos de acesso indicados & 13 & grupos operacionais. 16 endereços agrupados: APESP portal/item, BBM coleção/acervo geral e Wayback/Orkut; não instituições independentes. \href{https://github.com/murilloed/historia-riso-escrito-portugues/blob/main/data/base_registry.csv}{Dados no GitHub} \\
Fontes da seleção por acervo/década & 129 & identificadores de fonte. As 135 seleções reutilizam documentos entre quotas; 129 source\_id distintos. \href{https://github.com/murilloed/historia-riso-escrito-portugues/blob/main/data/evidence_register.csv}{Dados no GitHub} \\
Fontes selecionadas nos quatro inventários & 258 & identificadores de fonte. Reconciliação por source\_id e URL; não independência de obras, edições ou testemunhos. \href{https://github.com/murilloed/historia-riso-escrito-portugues/blob/main/data/consolidated_inventory.json}{Dados no GitHub} \\
Rubricas / marcadores miméticos e digitais do núcleo & 96 / 25 & registros. Partição do núcleo; rubricas descrevem reações, sem transcrever necessariamente o som. \href{https://github.com/murilloed/historia-riso-escrito-portugues/blob/main/data/consolidated_inventory.json}{Dados no GitHub} \\
Alternativas literárias de período & 2 & fontes. Bocage 1800 e Gonzaga 1812; não ocorrências nos exemplares fixos da Gazeta. \href{https://github.com/murilloed/historia-riso-escrito-portugues/blob/main/data/period_alternatives.csv}{Dados no GitHub} \\
Itens conferidos sem ocorrência localizada & 2 & itens. Leitura concluída; não prova de ausência na série nem quota positiva. \href{https://github.com/murilloed/historia-riso-escrito-portugues/blob/main/data/archive_decade_matrix.csv}{Dados no GitHub} \\
Confirmação no acervo indicado / fonte alternativa & 65 / 70 & linhas de quota. Partição das 135 seleções; confirmação alternativa não equivale a cotejo no portal original. \href{https://github.com/murilloed/historia-riso-escrito-portugues/blob/main/data/archive_decade_matrix.csv}{Dados no GitHub} \\
Classes lexicais / rubricas / miméticas da seleção & 112 / 18 / 5 & seleções. Partição das 135 linhas; não prevalência histórica nem fontes distintas. \href{https://github.com/murilloed/historia-riso-escrito-portugues/blob/main/data/consolidated_inventory.json}{Dados no GitHub} \\
Domínios registrados no manifesto & 87 & hosts. Inclui subdomínios, serviços técnicos e aquisições falhas; não equivale a bases consultadas integralmente. \href{https://github.com/murilloed/historia-riso-escrito-portugues/blob/main/data/source_domains.csv}{Dados no GitHub} \\
Fontes anteriores a 1900 & 40 & identificadores de fonte. 47 vínculos nos quatro inventários; inclui documentos desde 1645. Não é cobertura anual contínua. \href{https://github.com/murilloed/historia-riso-escrito-portugues/blob/main/data/consolidated_inventory.json}{Dados no GitHub} \\
\bottomrule\end{longtable}\normalsize

\small\begin{longtable}{>{\raggedright\arraybackslash}p{.29\textwidth}>{\raggedright\arraybackslash}p{.38\textwidth}>{\raggedright\arraybackslash}p{.22\textwidth}}
\caption{Inventário consolidado dos 16 endereços indicados: portais, coleções, item e captura. Sobreposições não são novas bases independentes.}\\
\toprule ID / recurso & Relação entre os recursos & Seleção e rastreabilidade\\\midrule\endfirsthead
\toprule ID / recurso & Relação entre os recursos & Seleção e rastreabilidade\\\midrule\endhead
H01 DocVirt / Jardinópolis \cite{r3archiveH01} & Portal/coleção; ver limitações de acesso e cobertura. & 11; referências por acervo e década \href{https://app.docvirt.com/dvd_ajhpmj/pageid/1}{Fonte}; \href{https://github.com/murilloed/historia-riso-escrito-portugues/blob/main/data/archive_decade_references.csv}{GitHub} \\
H02 O Globo \cite{r3archiveH02} & Portal/coleção; ver limitações de acesso e cobertura. & 7; referências por acervo e década \href{https://oglobo.globo.com/acervo/}{Fonte}; \href{https://github.com/murilloed/historia-riso-escrito-portugues/blob/main/data/archive_decade_references.csv}{GitHub} \\
H03 Folha \cite{r3archiveH03} & Portal/coleção; ver limitações de acesso e cobertura. & 11; referências por acervo e década \href{https://acervo.folha.com.br/index.do}{Fonte}; \href{https://github.com/murilloed/historia-riso-escrito-portugues/blob/main/data/archive_decade_references.csv}{GitHub} \\
H04 BORE / USP, Jornais \cite{r3archiveH04} & Portal/coleção; ver limitações de acesso e cobertura. & 13; referências por acervo e década \href{https://obrasraras.usp.br/xmlui/handle/123456789/4262}{Fonte}; \href{https://github.com/murilloed/historia-riso-escrito-portugues/blob/main/data/archive_decade_references.csv}{GitHub} \\
H05 Hemeroteca / BN \cite{r3archiveH05} & Portal/coleção; ver limitações de acesso e cobertura. & 21; referências por acervo e década \href{https://memoria.bn.gov.br/hdb/periodico.aspx}{Fonte}; \href{https://github.com/murilloed/historia-riso-escrito-portugues/blob/main/data/archive_decade_references.csv}{GitHub} \\
H06 Correio Braziliense \cite{r3archiveH06} & Portal/coleção; ver limitações de acesso e cobertura. & 6; referências por acervo e década \href{https://www.correiobraziliense.com.br/acervo}{Fonte}; \href{https://github.com/murilloed/historia-riso-escrito-portugues/blob/main/data/archive_decade_references.csv}{GitHub} \\
H07 Google News / Jornal do Brasil \cite{r3archiveH07} & Portal/coleção; ver limitações de acesso e cobertura. & 7; referências por acervo e década \href{https://news.google.com/newspapers?nid=0qX8s2k1IRwC}{Fonte}; \href{https://github.com/murilloed/historia-riso-escrito-portugues/blob/main/data/archive_decade_references.csv}{GitHub} \\
H08 APESP / portal \cite{r3archiveH08} & Portal APESP; H13 é um item deste mesmo acervo. & 0; referências por acervo e década \href{https://atom.arquivoestado.sp.gov.br/}{Fonte}; \href{https://github.com/murilloed/historia-riso-escrito-portugues/blob/main/data/archive_decade_references.csv}{GitHub} \\
H09 Senado / acervo de notícias \cite{r3archiveH09} & Portal/coleção; ver limitações de acesso e cobertura. & 3; referências por acervo e década \href{https://www12.senado.leg.br/noticias/acervo-historico}{Fonte}; \href{https://github.com/murilloed/historia-riso-escrito-portugues/blob/main/data/archive_decade_references.csv}{GitHub} \\
H10 Diario de Pernambuco / portal \cite{r3archiveH10} & Portal/coleção; ver limitações de acesso e cobertura. & 16; referências por acervo e década \href{https://www.diariodepernambuco.com.br/acervo}{Fonte}; \href{https://github.com/murilloed/historia-riso-escrito-portugues/blob/main/data/archive_decade_references.csv}{GitHub} \\
H11 CEPE / periódicos \cite{r3archiveH11} & Portal/coleção; ver limitações de acesso e cobertura. & 11; referências por acervo e década \href{https://acervo.cepe.com.br/}{Fonte}; \href{https://github.com/murilloed/historia-riso-escrito-portugues/blob/main/data/archive_decade_references.csv}{GitHub} \\
H12 BBM / Gazeta, coleção 8427 \cite{r3archiveH12} & Coleção Gazeta na BBM; H14 é o acervo geral da mesma biblioteca. & 0; referências por acervo e década \href{https://digital.bbm.usp.br/handle/bbm/8427}{Fonte}; \href{https://github.com/murilloed/historia-riso-escrito-portugues/blob/main/data/archive_decade_references.csv}{GitHub} \\
H13 APESP / item D000071 \cite{r3archiveH13} & Item individual do APESP; não é outra base independente. & 0; referências por acervo e década \href{https://atom.arquivoestado.sp.gov.br/br-spapesp-hemapesp-s001-b000565-d000071}{Fonte}; \href{https://github.com/murilloed/historia-riso-escrito-portugues/blob/main/data/archive_decade_references.csv}{GitHub} \\
H14 Brasiliana / BBM: acervo geral \cite{bbmAcervo2026} & Mesma biblioteca de H12; coleção geral, não outra instituição independente. & 14; títulos inéditos selecionados, rodada 4 \href{https://digital.bbm.usp.br/handle/bbm/1}{Fonte}; \href{https://github.com/murilloed/historia-riso-escrito-portugues/blob/main/data/revision4_new_titles.csv}{GitHub} \\
H15 Orkut preservado na Wayback \cite{r4orkutindex} & Conteúdo do Orkut preservado no serviço H16; não plataforma de preservação independente. & 1; comunidade selecionada, rodada 4 \href{https://web.archive.org/web/20141001005309/http://orkut.google.com/}{Fonte}; \href{https://github.com/murilloed/historia-riso-escrito-portugues/blob/main/data/revision4_new_titles.csv}{GitHub} \\
H16 Wayback Machine: portal geral \cite{waybackPortal2026} & Blog e fórum; Orkut contado separadamente em H15. O portal não foi pesquisado exaustivamente. & 2; blog e fórum selecionados, rodada 4 \href{https://web.archive.org/}{Fonte}; \href{https://github.com/murilloed/historia-riso-escrito-portugues/blob/main/data/revision4_new_titles.csv}{GitHub} \\
\bottomrule\end{longtable}\normalsize


Um zero na coluna de seleção indica ausência de referência autenticada naquele recorte e rodada, não ausência de consulta ao portal nem inexistência histórica de riso.

Os dados completos estão em \href{https://github.com/murilloed/historia-riso-escrito-portugues/blob/main/data/base_registry.csv}{bases}, \href{https://github.com/murilloed/historia-riso-escrito-portugues/blob/main/data/method_quantification.csv}{quantificação}, \href{https://github.com/murilloed/historia-riso-escrito-portugues/blob/main/data/source_domains.csv}{domínios} e \href{https://github.com/murilloed/historia-riso-escrito-portugues/blob/main/data/downloads.json}{aquisições}. Os links levam à versão corrente; a reprodução deve identificar a versão ou o commit consultado. O inventário explicita pendências e não equivale à cobertura integral das coleções.

\section{Fundamentação estatística da quantidade de documentos por ano}
\subsection{Por que três referências não constituem justificativa estatística}
O roteiro inicial pede três jornais e três livros por período. Esse número é uma quota editorial útil para organizar exemplos, mas não fornece representatividade por si só. A existência de três títulos conhecidos não informa o número de páginas examinadas, a quantidade de oportunidades de riso, o universo de documentos preservados ou a dependência entre passagens. A repetição de exemplos não estima frequência social. Tampouco um cálculo de amostra transforma retrospectivamente a busca dirigida em sorteio probabilístico.

É preciso separar dois objetivos. Para demonstrar \emph{existência}, uma atestação autêntica basta para aquele documento e data. Para estimar \emph{prevalência} ou comparar anos, é necessário definir população, unidade amostral, mecanismo de seleção e incerteza. Para demonstrar \emph{origem}, nenhum tamanho de amostra finito elimina automaticamente a possibilidade de documentos anteriores perdidos ou não inspecionados. A estatística auxilia o dimensionamento e a descrição; não produz uma certidão de nascimento etimológico.

\subsection{Desenho prospectivo estratificado por ano e suporte}
Propõe-se construir uma moldura de documentos para cada ano $y$ e suporte $s$, distinguindo ao menos livros e periódicos. O estrato é $h=(y,s)$ e contém $N_h$ unidades elegíveis. Livros podem ser definidos por edição publicada naquele ano; periódicos, por número, com páginas selecionadas em segundo estágio quando necessário. Interações digitais formam estratos próprios, pois sessões, mensagens e usuários não são equivalentes a volumes impressos. País, gênero e acervo devem ser registrados para controlar composição.

O estimando básico sugerido é a proporção de documentos do estrato com pelo menos uma representação gráfica de riso: $p_h=P(I_h=1)$. Outro estimando possível é a frequência de determinado tipo por quantidade de palavras inspecionadas. Esses objetos não se confundem. Uma obra longa tende a oferecer mais oportunidades de ocorrência, razão pela qual também se registra extensão. A frequência normalizada, se houver denominadores confiáveis, será $f_h=10^6 C_h/W_h$, em que $C_h$ é número de ocorrências e $W_h$ o total de palavras elegíveis.

Quando os estratos têm dimensões conhecidas, pode-se fazer censo dos pequenos e sorteio dos grandes, registrando probabilidades de inclusão. A estimativa agregada usa pesos proporcionais aos universos, e não uma média simples de estratos com tamanhos diferentes. Se custos e variâncias forem conhecidos, a alocação ótima é proporcional a $N_hS_h/\sqrt{c_h}$, em que $S_h$ é o desvio-padrão do indicador e $c_h$ o custo por unidade. Sem piloto ou moldura, essa alocação é uma proposta, não um procedimento executado \cite{cochran1977}.

\subsection{Dimensionamento para uma proporção}
Sob amostragem aleatória simples, população grande, nível nominal de 95\% e erro absoluto $e$, um dimensionamento inicial é
\begin{equation}
n_0=\frac{z_{0,975}^2p(1-p)}{e^2},\qquad z_{0,975}=1{,}96.
\end{equation}
Na ausência de proporção piloto, $p=0{,}5$ maximiza a variância. Com $e=0{,}10$, resulta $n_0=96{,}04$, arredondado para 97 documentos \emph{por estrato} se a precisão for exigida em cada combinação de ano e suporte. Para $e=0{,}05$, são 385. Esses números não significam 97 ou 385 referências bibliográficas a citar no artigo: as referências identificam os documentos examinados, enquanto tabelas exemplificativas podem exibir apenas uma seleção transparente.

Para universo finito conhecido, a correção é
\begin{equation}
n_h=\left\lceil\frac{N_hn_0}{N_h+n_0-1}\right\rceil.
\end{equation}
Utiliza-se o valor não arredondado de $n_0$ na fórmula. Se $N_h=20$, o resultado para $e=0{,}10$ é 17; se $N_h=100$, é 50; se $N_h=1000$, é 88. Pressupõe-se sorteio sem reposição e população adequadamente coberta. A precisão nominal não corrige vieses de preservação, OCR ou seleção do acervo.

\subsection{Dimensionamento para detecção de uma forma rara}
Uma pergunta mais adequada à busca de atestações é: qual a chance de encontrar pelo menos um documento positivo? Sob ensaios independentes com probabilidade constante $p$, tem-se
\begin{equation}
P(\text{detecção}\mid n,p)=1-(1-p)^n,
\qquad n\geq\left\lceil\frac{\log(1-q)}{\log(1-p)}\right\rceil,
\end{equation}
em que $q$ é a probabilidade de detecção desejada. Para $q=0{,}95$ e $p=0{,}05$, são necessários 59 documentos; para $p=0{,}01$, 299. Com três documentos, a chance é apenas 14,3\% no primeiro cenário e 3,0\% no segundo. Seis documentos elevam esses valores a 26,5\% e 5,9\%. A fórmula explicita por que uma quota pequena não fundamenta afirmações de ausência.

\begin{table}[htbp]\centering\small
\caption{Probabilidade de ao menos uma detecção em cenários hipotéticos independentes. Não são resultados observados.}
\begin{tabular}{rrr}\toprule
Documentos ($n$) & $p=0{,}05$ & $p=0{,}01$ \\\midrule
3 & 14,3\% & 3,0\% \\
6 & 26,5\% & 5,9\% \\
30 & 78,5\% & 26,0\% \\
59 & 95,2\% & 44,7\% \\
97 & 99,3\% & 62,3\% \\
299 & 100,0\% & 95,0\% \\
\bottomrule\end{tabular}
\end{table}

A independência é uma aproximação. Notícias reproduzidas, obras do mesmo autor e sessões com os mesmos participantes podem aumentar a correlação. Se necessário, um planejamento por conglomerados deverá considerar efeito de desenho $D$, com ajuste inicial $n_{\mathrm{aj}}=\lceil D n_0\rceil$, seguido de análise compatível com a amostragem efetiva. Não se adota um valor de $D$ como se fosse estimado pelos dados atuais.

\subsection{Incerteza, zeros e análise temporal}
Para uma amostra aleatória simples binomial, o intervalo de Wilson é uma alternativa ao intervalo normal simétrico, particularmente com contagens pequenas \cite{wilson1927}. Com $\hat p=x/n$, seus limites são
\begin{equation}
\frac{\hat p+z^2/(2n)\ \pm\ z\sqrt{\hat p(1-\hat p)/n+z^2/(4n^2)}}{1+z^2/n}.
\end{equation}
O script inclui a fórmula para reprodução do planejamento. Ela não é aplicada como intervalo populacional à coleta por conveniência deste artigo. Quando não há amostra válida num ano, informa-se ``sem atestação confirmada nesta coleta'', e não uma frequência zero da língua. O denominador precisa ser exposto mesmo quando a contagem de positivos é zero.

Com corpus ampliado, comparações de contagens poderão utilizar modelos com compensação pela extensão textual, efeitos de documento e controles de gênero, país e acervo. Modelos de mudança temporal não devem ser ajustados a uma tabela composta apenas de exemplos selecionados. A forma mais antiga encontrada também é uma estatística dependente da cobertura: um acervo maior pode deslocá-la para trás sem que a história real tenha mudado.

\subsection{Cobertura anual efetiva e plano de expansão}
A rodada de 8 e 9 de novembro de 2026 buscou um testemunho adicional por ano ainda descoberto. Foram incluídos \AnnualAdded{} documentos, cobrindo \AnnualCovered{} dos 110 anos de 1900 a 2009. A seleção dirigida priorizou os Anais do Senado Federal brasileiro, os diários parlamentares portugueses e periódicos brasileiros da Hemeroteca Digital. Entre os novos testemunhos, 87 são registros parlamentares e nove são periódicos. Há 73 registros brasileiros e 23 portugueses. As rubricas editoriais registram o riso numa situação narrada ou transcrita. As grafias miméticas representam o som por escrito. Essas categorias foram mantidas distintas. O OCR serviu à localização em fac-símiles, com conferência da passagem na página. As rubricas eletrônicas foram cotejadas com o contexto na transcrição institucional. O ano segue o documento identificado pelo cabeçalho, edição ou catálogo, conforme anotado no localizador. O protocolo visou comprovar ao menos uma ocorrência por ano e privilegiou séries pesquisáveis. A seleção, portanto, concentra gêneros e instituições. Os dados estão disponíveis na \href{https://github.com/murilloed/historia-riso-escrito-portugues/blob/main/data/annual_coverage.csv}{série anual} e no \href{https://github.com/murilloed/historia-riso-escrito-portugues/blob/main/data/annual_supplement.csv}{inventário dos novos testemunhos}, com URL, forma, nível de evidência e localizador.

O planejamento de 97 documentos por estrato seria muito mais extenso que a quota inicial. Para 110 anos e dois suportes, corresponderia a até 21.340 documentos antes de correções por universo finito, indisponibilidade ou conglomerados. Essa ordem de grandeza é uma consequência da precisão escolhida, não uma exigência universal. Um projeto pode adotar precisão menor, estratos decenais ou objetivo apenas de detecção, desde que justifique a decisão e preserve os limites inferenciais.

\section{Resultados documentais}
\subsection{Testemunhos confirmados e suas datas}
A Tabela~\ref{tab:attestations} reúne passagens selecionadas com localizadores. Ela não é uma série de frequências e não representa a proporção de brasileiros que empregava cada grafia. As datas literárias identificam a obra; quando o acesso é uma transcrição posterior, isso é explicitado. As datas de estudos acadêmicos não são convertidas automaticamente em datas das interações.

\small
\begin{longtable}{p{0.07\textwidth}p{0.24\textwidth}p{0.17\textwidth}p{0.37\textwidth}}
\caption{Atestações selecionadas: ano de referência, grafia e localizador. E1: imagem; E2: transcrição primária; E3: interação reportada.}\label{tab:attestations}\\
\toprule Ano & Documento & Grafia & Localizador / evidência \\\midrule\endfirsthead
\toprule Ano & Documento & Grafia & Localizador / evidência \\\midrule\endhead
1872 & Til \cite{alencar1872} & Quiá!... quiá!... quiá!... & Volume II, capítulo VIII; âncora: Ai, menina! (E2). \\
1879 & A Imprensa Ytuana : jornal scientifico, litterario, noticiozo e industrial do Instituto do Novo-Mundo, ano 004, n. 183, 1879 \cite{r3doc1} & (Riso) & 1879-09-13; PDF p. 2 (E1). \\
1884 & A segunda vida \cite{assis1884} & cá, cá, cá, cá, cá... & Âncora: Era uma gargalhada assim (E2). \\
1900 & Anais do Senado Federal, livro 1 \cite{anual1900} & (Riso.) & PDF p. 248 (E2). \\
1901 & A cidade e as serras \cite{queiros1901} & Ah! ah! ah! & Âncora: Então tambem eu rio!; PG 18220 (E2). \\
1902 & Anais do Senado Federal, livro 1 \cite{anual1902} & (Riso) & PDF p. 60 (E2). \\
1903 & Anais do Senado Federal, livro 1 \cite{anual1903} & (Riso.) & PDF p. 114 (E2). \\
1904 & Anais do Senado Federal, livro 1 \cite{anual1904} & (Riso.) & PDF p. 320 (E2). \\
1905 & Anais do Senado Federal, livro 1 \cite{anual1905} & (riso) & PDF p. 88 (E2). \\
1906 & Anais do Senado Federal, livro 1 \cite{anual1906} & (Risos) & PDF p. 318 (E2). \\
1907 & Anais do Senado Federal, livro 1 \cite{anual1907} & (Riso.) & PDF p. 225 (E2). \\
1908 & Anais do Senado Federal, livro 3 \cite{anual1908} & (Riso.) & PDF p. 169 (E2). \\
1909 & Anais do Senado Federal, livro 2 \cite{anual1909} & (Riso.) & PDF p. 21 (E2). \\
1910 & Anais do Senado Federal, livro 2 \cite{anual1910} & (Riso) & PDF p. 153 (E2). \\
1911 & Anais do Senado Federal, livro 1 \cite{anual1911} & (Riso nas galerias.) & PDF p. 312 (E2). \\
1912 & Anais do Senado Federal, livro 1 \cite{anual1912} & (Risos.) & PDF p. 418 (E2). \\
1913 & Anais do Senado Federal, livro 5 \cite{anual1913} & (Risos nas galerias.) & p. 238, PDF p. 255 (E2). \\
1914 & Anais do Senado Federal, livro 1 \cite{anual1914} & (Risos.) & PDF p. 205 (E2). \\
1915 & Anais do Senado Federal, livro 1 \cite{anual1915} & (riso) & PDF p. 121 (E2). \\
1916 & Anais do Senado Federal, livro 4 \cite{anual1916} & (Risos.) & sessão de 24/08/1916, p. 399, PDF p. 415 (E2). \\
1917 & A engomadeira \cite{almada1917} & Ah! Ah! Ah!... & Âncora: E 'inda o confessa!; PG 23879 (E2). \\
1918 & Diário da Câmara dos Deputados, n. 008 \cite{anual1918} & Risos. & data no catálogo 1918-08-02, p. 18 (E2). \\
1919 & Urupês, quinta edição \cite{lobato1919} & Quá, quá, quá! & p. impressa 19; PDF p. 24; conto O engraçado arrependido (E1). \\
1920 & Diário da Câmara dos Deputados, n. 059 \cite{anual1920} & Risos. & data no catálogo 1920-03-31, p. 41 (E2). \\
1921 & Anais do Senado Federal, livro 1 \cite{anual1921} & (Riso.) & PDF p. 500 (E2). \\
1922 & Anais do Senado Federal, livro 3 \cite{anual1922} & (Riso.) & PDF p. 12 (E2). \\
1923 & Anais do Senado Federal, livro 2 \cite{anual1923} & (Riso.) & PDF p. 92 (E2). \\
1924 & Anais do Senado Federal, livro 1 \cite{anual1924} & (Riso.) & PDF p. 611 (E2). \\
1925 & Diário da Câmara dos Deputados, n. 77 \cite{portugal1925} & Risos. & 11/06/1925; p. 14; intervenção de Cunha Leal (E2). \\
1926 & Anais do Senado Federal, livro 1 \cite{anual1926} & (Risos.) & PDF p. 311 (E2). \\
1927 & Fon-Fon \cite{anual1927} & Quá, quá & PDF p. 75; Saibam todos (E1). \\
1928 & Diário Nacional, n. 290 \cite{diarionacional1928} & (risos, muito bem) & 17/06/1928; p. 4; discurso de Mattos Pimenta (E1). \\
1929 & O Tico-Tico \cite{anual1929} & Quá, quá, quá, quá.... & 12/06/1929; PDF p. 16; A raposa e o gallo; riso ficcional explícito (E1). \\
1930 & A Batalha \cite{anual1930} & Quá-quá-quá...! & 06/08/1930, PDF p. 2, seção A pedido (E1). \\
1931 & O Tico-Tico \cite{anual1931} & quá quá, quá... & 22/04/1931; PDF p. 23; Ri... quá quá, quá... (E1). \\
1932 & Diario de Noticias \cite{anual1932} & Quá-quá! & 02/10/1932; PDF p. 12; comentário de xadrez assinado Pharaó (E1). \\
1933 & Anais da Assembleia Nacional Constituinte, livro 1 \cite{anual1933} & (Riso.) & PDF p. 318 (E2). \\
1934 & Anais da Assembleia Nacional Constituinte, livro 6 \cite{anual1934} & (Risos.) & PDF p. 79 (E2). \\
1935 & O Tico-Tico, n. 1569 \cite{ticotico1935} & Quá quá quá... & 30/10/1935; capa, último quadro; PDF p. 1 (E1). \\
1936 & Diário das Sessões da Assembleia Nacional, n. 075 \cite{anual1936} & (Risos) & sessão 1936-02-07; p. 437 (E2). \\
1937 & Diário das Sessões da Assembleia Nacional, n. 117 \cite{anual1937} & (Risos) & sessão 1937-02-11; p. 443 (E2). \\
1938 & O Tico-Tico \cite{anual1938} & Ha ha ha ha! & 01/06/1938; PDF p. 18; As aventuras do Camondongo Mickey (E1). \\
1939 & Diário das Sessões da Assembleia Nacional, n. 053 \cite{anual1939} & (Risos) & sessão 1939-12-11; p. 33 (E2). \\
1940 & O Tico-Tico \cite{anual1940} & Quá - quá! quá! & 22/05/1940; PDF p. 26; As aventuras do Camondongo Mickey (E1). \\
1941 & Diário Carioca, n. 3902 \cite{carioca1941} & Ah! ah! & 08/03/1941; p. 13; folhetim As memórias de Sherlock Holmes (E1). \\
1942 & Diário das Sessões da Assembleia Nacional, n. 008 \cite{anual1942} & Risos. & data no catálogo 1942-12-12, p. 49 (E2). \\
1943 & Diário das Sessões da Assembleia Nacional, n. 024 \cite{anual1943} & Risos. & data no catálogo 1943-03-17, p. 235 (E2). \\
1944 & Diário das Sessões da Assembleia Nacional, n. 066 \cite{anual1944} & (Risos) & sessão 1944-03-22; p. 298 (E2). \\
1945 & Diário das Sessões da Assembleia Nacional, n. 118 \cite{anual1945} & (risos) & sessão 1945-01-11; p. 115 (E2). \\
1946 & Anais do Senado Federal, livro 1 \cite{anual1946} & (Risos.) & PDF p. 76 (E2). \\
1947 & Jornal de Notícias, n. 267 \cite{jornalnoticias1947} & (Risos dos presentes) & 02/03/1947; p. 1; fala de Bidault; não é Folha da Manhã (E1). \\
1948 & Anais do Senado Federal, livro 2 \cite{anual1948} & (Risos) & PDF p. 501 (E2). \\
1949 & Diário das Sessões da Assembleia Nacional, n. 180 \cite{anual1949} & Risos. & sessão 1949-03-24; p. 325 (E2). \\
1950 & Carioca, n. 792 \cite{anual1950} & Quá, quá, quá & 07/12/1950, PDF p. 48, coluna Por trás do dial (E1). \\
1951 & Diário das Sessões da Assembleia Nacional, n. 105 \cite{anual1951} & Risos. & data no catálogo 1951-04-28, p. 984 (E2). \\
1952 & Diário das Sessões da Assembleia Nacional, n. 132 \cite{anual1952} & (risos) & data no catálogo 1952-03-04, p. 427 (E2). \\
1953 & Anais do Senado Federal, livro 1 \cite{anual1953} & (Risos). & PDF p. 14 (E2). \\
1954 & Anais do Senado Federal, livro 1 \cite{anual1954} & (Riso). & PDF p. 7 (E2). \\
1955 & Diário do Congresso Nacional, suplemento \cite{camara1955} & (Risos) & 10/12/1955; p. 13; intervenção de Alberto Torres (E1). \\
1956 & Mundo Esportivo \cite{anual1956} & quá, quá, quá & 25/05/1956, PDF p. 2, diálogo humorístico (E1). \\
1957 & Anais do Senado Federal, livro 1 \cite{anual1957} & (Risos). & PDF p. 230 (E2). \\
1958 & Anais do Senado Federal, livro 1 \cite{anual1958} & (Riso). & PDF p. 662 (E2). \\
1959 & Anais do Senado Federal, livro 2 \cite{anual1959} & (Riso) & PDF p. 398 (E2). \\
1960 & Anais do Senado Federal, livro 2 \cite{anual1960} & (Riso) & PDF p. 219 (E2). \\
1961 & Anais do Senado Federal, livro 1 \cite{anual1961} & (Riso). & PDF p. 159 (E2). \\
1962 & Anais do Senado Federal, livro 1 \cite{anual1962} & (Riso.) & PDF p. 190 (E2). \\
1963 & Diário das Sessões da Assembleia Nacional, n. 066 \cite{anual1963} & (Risos) & data no catálogo 1963-01-08, p. 1747 (E2). \\
1964 & Anais do Senado Federal, livro 2 \cite{anual1964} & (Riso) & PDF p. 25 (E2). \\
1965 & Anais do Senado, livro 1 \cite{senado1965} & (Risos.) & p. impressa 162; PDF p. 185; intervenção de Edmundo Levi (E1). \\
1966 & Anais do Senado Federal, livro 1 \cite{anual1966} & (Riso) & PDF p. 148 (E2). \\
1967 & Diário das Sessões da Assembleia Nacional, n. 059 \cite{anual1967} & (Risos) & sessão 1967-01-19; p. 1059 (E2). \\
1968 & Anais do Senado Federal, livro 1 \cite{anual1968} & (Risos.) & PDF p. 137 (E2). \\
1969 & Anais do Senado Federal, livro 1 \cite{anual1969} & (Riso.) & PDF p. 164 (E2). \\
1970 & Anais do Senado Federal, livro 4 \cite{anual1970} & (Risos.) & p. 215, PDF p. 216 (E2). \\
1971 & Anais do Senado Federal, livro 1 \cite{anual1971} & (Risos.) & PDF p. 233 (E2). \\
1972 & Anais do Senado Federal, livro 1 \cite{anual1972} & (risos) & PDF p. 68 (E2). \\
1973 & Anais do Senado Federal, livro 1 \cite{anual1973} & (Risos.) & PDF p. 124 (E2). \\
1974 & Anais do Senado Federal, livro 1 \cite{anual1974} & (Risos.) & PDF p. 344 (E2). \\
1975 & Anais do Senado Federal, livro 1 \cite{anual1975} & (risos) & PDF p. 173 (E2). \\
1976 & Anais do Senado Federal, livro 1 \cite{anual1976} & (Risos.) & PDF p. 49 (E2). \\
1977 & Anais do Senado Federal, livro 1 \cite{anual1977} & (Risos.) & PDF p. 83 (E2). \\
1978 & Anais do Senado Federal, livro 1 \cite{anual1978} & (Risos.) & PDF p. 56 (E2). \\
1979 & Correio do Povo \cite{correio1979} & AH, AH, AH & 28/04/1979; p. 5; tira Bidu, quadros finais (E1). \\
1980 & Anais do Senado Federal, livro 1 \cite{anual1980} & (Risos.) & PDF p. 35 (E2). \\
1981 & Anais do Senado Federal, livro 2 \cite{anual1981} & (risos) & PDF p. 81 (E2). \\
1982 & Anais do Senado Federal, livro 1 \cite{anual1982} & (Risos.) & PDF p. 100 (E2). \\
1983 & Anais do Senado Federal, livro 1 \cite{anual1983} & (Risos.) & PDF p. 57 (E2). \\
1984 & Anais do Senado Federal, livro 3 \cite{anual1984} & (Risos.) & PDF p. 281 (E2). \\
1985 & Diário do Congresso Nacional, seção I \cite{camara1985} & (risos e aplausos.) & 29/03/1985; p. impressa 2107; PDF p. 55 (E1). \\
1986 & Anais do Senado Federal, livro 3 \cite{anual1986} & (Risos.) & PDF p. 13 (E2). \\
1987 & Anais do Senado Federal, livro 1 \cite{anual1987} & (Risos.) & PDF p. 129 (E2). \\
1988 & Anais do Senado Federal, livro 1 \cite{anual1988} & (risos.) & PDF p. 5 (E2). \\
1989 & Anais do Senado Federal, livro 1 \cite{anual1989} & (Risos) & PDF p. 82 (E2). \\
1990 & Anais do Senado Federal, livro 1 \cite{anual1990} & (Risos) & PDF p. 48 (E2). \\
1991 & Anais do Senado Federal, livro 1 \cite{anual1991} & (Risos.) & PDF p. 43 (E2). \\
1992 & Anais do Senado Federal, livro 2 \cite{anual1992} & (Risos) & diário de 20/02/1992, p. 358, PDF p. 110 (E2). \\
1993 & Diário da Assembleia da República, I série, n. 073 \cite{anual1993} & Risos do PSD e do CDS. & data no catálogo 1993-05-20, p. 2300 (E2). \\
1994 & Folha de S.Paulo \cite{folha1994} & (risos) & 03/09/1994; Veja íntegra da conversa de Ricupero com Monforte (E2). \\
1995 & Anais do Senado Federal, livro 1 \cite{anual1995} & (Risos) & PDF p. 144 (E2). \\
1996 & Diário da Assembleia da República, I série, n. 032 \cite{anual1996} & Risos. & data no catálogo 1996-01-26, p. 911 (E2). \\
1997 & Anais do Senado Federal, livro 2 \cite{anual1997} & (risos) & janeiro de 1997, entrevista reproduzida nos Anais, p. 13, PDF p. 39 (E2). \\
1998 & Diário da Assembleia da República, I série, n. 002 \cite{anual1998} & Risos do PSD. & data no catálogo 1998-09-17, p. 40 (E2). \\
1999 & Anais do Senado Federal, livro 1 \cite{anual1999} & (Risos) & PDF p. 507 (E2). \\
2000 & Diário da Assembleia da República, I série, n. 033 \cite{anual2000} & Risos do PCP. & data no catálogo 2000-12-21, p. 1292 (E2). \\
2001 & Diário da Assembleia da República, I série, n. 034 \cite{anual2001} & Risos do CDS-PP e do PSD. & data no catálogo 2001-01-03, p. 1372 (E2). \\
2002 & Diário da Assembleia da República, I série, n. 031 \cite{anual2002} & Risos. & data no catálogo 2002-01-23, p. 1205 (E2). \\
2003 & Diário da Assembleia da República, I série, n. 071 \cite{anual2003} & Risos do PCP. & data no catálogo 2003-01-08, p. 3008 (E2). \\
2004 & Diário da Assembleia da República, I série, n. 035 \cite{anual2004} & Risos do PCP. & data no catálogo 2004-01-07, p. 2001 (E2). \\
2005 & Anais do Senado Federal, livro 1 \cite{anual2005} & (Risos.) & PDF p. 606 (E2). \\
2006 & Anais do Senado Federal, livro 1 \cite{anual2006} & (Risos.) & PDF p. 361 (E2). \\
2007 & Chat reproduzido por Young (2008) \cite{young2008} & kkkk & Dissertação p. 93; contexto: semestre 2007.2 (E3). \\
2008 & Anais do Senado Federal, livro 1 \cite{anual2008} & (Risos.) & PDF p. 957 (E2). \\
2009 & Diário da Assembleia da República, I série, n. 030 \cite{anual2009} & Risos do PS. & data no catálogo 2009-01-07, p. 49 (E2). \\
2010 & Leão (2010) \cite{leao2010} & kkkkkkkk / rsrsrsrs & PDF pp. 90-92; fragmentos 97, 101 e 103 e análise (E3). \\
2010 & Cifra Club: como afinar meio tom abaixo \cite{r4forumcifraclub} & kkkkkkkkkkkkk & resposta Enviado: 12/jul/10 15:21; após boa!; captura 2010-09-15T13:54:14Z (E2). \\
2014 & Orkut: comunidade Kiko Dança Psy-TRANCE \cite{r4orkutcomunidades} & hehehe & descrição da comunidade, após sempre agitado; captura 2014-10-04T09:58:23Z (E2). \\
2015 & Silva (2015) \cite{silva2015} & Kkkk & p. impressa 86, PDF p. 87; CA7 P2 (E3). \\
2016 & A Cidade de Jardinópolis \cite{r3doc18} & Hahaha... & 2016-01-09; pasta 49, página lógica 3; p. impressa 3; Só Rindo, Claudio Aguilhera (E1). \\
2019 & Macunaíma, edição UFFS \cite{andrade2019} & Quá! quá! quá quaquá!... & p. impressa e PDF 139; ficha catalográfica PDF 205 (E1). \\
2022 & A Cidade de Jardinópolis \cite{r3doc19} & kkkkk & 2022-01-29; pasta 57, página lógica 21; p. impressa 5; Entrevista (E1). \\
2025 & Correio Braziliense: Uma sintonia capital \cite{r3doc21} & (risos) & 2025-04-07; PDF p.22; Diversão \& Arte; entrevista de Walter Lima Jr. por Ricardo Daehn (E1). \\
\bottomrule\end{longtable}\normalsize

Em \textit{Til}, segundo volume, capítulo VIII, a reação de Nhá Tudinha contém ``Quiá!... quiá!... quiá!...'' \cite{alencar1872}. O contexto descreve uma risada e permite classificar a sequência como representação de riso. A passagem confirma a grafia no texto consultado, mas não estabelece que seja a primeira onomatopeia de riso em português. Também não contém \forma{kkkk}. A autoria é de José de Alencar; atribuí-la genericamente a Eça é um erro de identificação.

Em \textit{A segunda vida}, de Machado de Assis, o corpus eletrônico consultado apresenta ``cá, cá, cá, cá, cá...'' após a descrição de uma gargalhada \cite{assis1884}. O caso mostra a relação entre encenação literária e representação sonora. A presença do mesmo mecanismo repetitivo não revela por si só se o autor conhecia determinada passagem de outro escritor ou se um usuário de chat a herdou mais tarde.

Em \textit{A cidade e as serras}, publicado em 1901, encontra-se uma sequência de \forma{ah} associada explicitamente a rir \cite{queiros1901}. A edição eletrônica permite localizar a passagem pela âncora ``Então tambem eu rio!'' ou sua variante efetiva indicada no inventário. Em \textit{A engomadeira}, de 1917, a sequência ``Ah! Ah! Ah!...'' integra uma cena de zombaria \cite{almada1917}. Esses exemplos impedem reduzir a história do riso a uma sucessão em que um formato substitui integralmente outro a cada década.

Na capa de \textit{O Tico-Tico}, número 1569, de 30 de outubro de 1935, a narrativa gráfica ``A caricatura do Bolinha'' apresenta \forma{quá quá quá} no último quadro \cite{ticotico1935}. A imagem foi inspecionada para resolver ruídos do OCR. O testemunho é particularmente útil porque reúne data, número, localização e contexto visual. Demonstra uso numa revista infantil brasileira antes da cronologia que atribuía sua consolidação inicial aos quadrinhos Disney dos anos 1950 ou 1960; não determina, isoladamente, o alcance social da forma.


O \textit{Diário Carioca}, em 8 de março de 1941, publica uma sequência de \forma{ah} no folhetim traduzido de Conan Doyle, página 13 \cite{carioca1941}. A data identifica a publicação brasileira da tradução, não a composição original. No \textit{Correio do Povo} de 28 de abril de 1979, página 5, a tira do Bidu contém repetições de \forma{AH} nos quadros finais \cite{correio1979}. A imagem permite corrigir o OCR, que alternava equivocadamente \forma{AU} e \forma{AH}. O caso demonstra permanência dessa convenção em quadrinhos de 1979, sem medir frequência no conjunto da imprensa.

\subsection{A documentação digital de 2007 e os estudos posteriores}
\subsubsection{Rubricas institucionais e jornalísticas: 1925 a 1994}
As rubricas foram incluídas em categoria própria porque representam uma reação descrita, não uma imitação fonética. No Diário da Câmara dos Deputados de Portugal, sessão de 11 de junho de 1925, página 14, a intervenção de Cunha Leal é seguida de \forma{Risos.} na transcrição oficial \cite{portugal1925}. O portal informa que a extração é automática; o testemunho foi classificado como E2, sem equiparação a colação do fac-símile.

No suplemento do Diário do Congresso Nacional de 10 de dezembro de 1955, página 13, a intervenção de Alberto Torres recebe a marca \forma{(Risos)} \cite{camara1955}. Nos Anais do Senado de 1965, livro 1, página impressa 162 e página 185 do PDF, a fala de Edmundo Levi sobre a liberdade dos senadores e um cafezinho é acompanhada de \forma{(Risos.)} \cite{senado1965}. A imagem confirma o sinal; o dia da sessão não foi individualizado neste registro, preservando-se apenas o ano do volume.

No Diário do Congresso Nacional de 29 de março de 1985, página 2107, a marca \forma{(risos e aplausos.)} aparece em material transcrito inserido nos trabalhos \cite{camara1985}. Na versão eletrônica da Folha de S.Paulo de 3 de setembro de 1994, a transcrição da conversa de Ricupero com Monforte inclui \forma{(risos)} após uma anedota \cite{folha1994}. Esses dados cobrem modalidades descritivas em diferentes décadas; não autorizam calcular predominância nem equiparar rubrica e onomatopeia.

Com as novas atestações, o inventário tem ao menos um documento de cada década de 1900 a 2000. Essa cobertura decenal mínima não significa cobertura de todos os anos, três livros e três periódicos por período, nem amostra equilibrada entre suportes. As tabelas anuais explicitam a diferença.

\subsubsection{Interações digitais reproduzidas em pesquisas}
Young \cite{young2008} reproduz, na página 93, um diálogo de bate-papo com cadeias de \forma{k}; a pesquisa é situada no semestre 2007.2. O registro é um testemunho reportado de interação digital, acessível em fonte acadêmica de 2008. O dado autoriza afirmar que \forma{kkkk} é documentado naquele conjunto ligado a 2007. Não autoriza concluir que surgiu em 2007 ou que aquela plataforma tenha sido seu berço nacional. O download direto do PDF falhou no ambiente local, embora o documento e o registro institucional tenham sido consultados pelo serviço de leitura web; essa diferença consta no inventário.

Leão \cite{leao2010} reúne, em pesquisa publicada em 2010, exemplos de chat com \forma{kkkk}, \forma{rsrsrsrs} e sequências heterogêneas, nas páginas de arquivo 90--92. Silva \cite{silva2015} registra \forma{Kkkk} no comentário CA7 P2, página impressa 86, correspondente à página 87 do PDF, e documenta outras formas nos anexos. Esses casos mostram coexistência de repertórios em contextos educacionais. Não constituem prova de que um tipo tenha desaparecido nacionalmente quando outro passou a circular.

O ano de 2010 ou 2015 é o ano do documento consultado quando não há data individual validada da mensagem. Preservar essa distinção evita construir uma linha temporal com precisão inexistente. Do mesmo modo, um artigo recente que reproduz comentários antigos deve fornecer metadados da coleta para que os comentários sejam datados independentemente da publicação.

\subsection{Busca, downloads e falsos positivos}
O conjunto baixado inclui textos literários sem sequência de riso confirmada pela busca dirigida, PDFs de periódicos e documentos de apoio. Entre os resultados de busca, ``Quá... Quá...'' aparece também como nome comercial em \textit{Jornal de Notícias}, de 1946, e como parte do título de um programa na revista \textit{Intervalo}, de 1965. Esses usos podem ser pertinentes à circulação cultural de uma expressão, mas não são contados como transcrições de alguém rindo.

As buscas em três números de \textit{O Malho}, de 1920, não produziram atestação que tivesse sido validada para inclusão na tabela principal. Uma correspondência automática atravessava a palavra \forma{ahi} e a forma verbal \forma{ha}, gerando um falso positivo. O resultado demonstra a necessidade de fronteiras de palavra e conferência contextual. Não demonstra ausência de riso onomatopaico em todo o periódico ou na década.

\subsection{A auditoria da série de 1900 a 2000}
O material inicial apresenta onze blocos decenais, cada um com exemplo literário e jornalístico supostamente literal. Sua apresentação dá aparência de uma série documental contínua, mas faltam páginas, edições específicas e, em vários casos, identificação suficiente de autor, seção ou número. A matriz no apêndice conserva os 22 alvos de conferência. Nenhuma dessas frases é promovida a citação autêntica apenas pela existência da obra ou do periódico citado.

O problema não é apenas bibliográfico. A série associa 1900 a solenidade, 1910 a caricatura, 1920 a riso de salão, 1930 a escárnio, 1940 a rubricas, 1950 a Disney, 1960 a contracultura, 1970 a gibis, 1980 a datilografia, 1990 a \forma{rs} e 2000 a \forma{kkkk}. Sem contagens comparáveis, trata-se de uma narrativa interpretativa, não de uma distribuição observada. A existência de rubricas institucionais ou de quadrinhos numa época não implica predomínio nacional de uma grafia.

As alegações de que a norma exigia sempre \forma{h} final, de que ninguém agrupava letras em 1900 e de que o \forma{rs} nasceu para economizar fita datilográfica não receberam documentação normativa ou material suficiente. A explicação econômica é testável, mas não constitui prova em si mesma. Os exemplos inventados para ilustrar um estilo, quando apresentados como transcrição, devem ser removidos da condição de fonte histórica.

\section{Ampliação por acervo e década: resultados e limites}
\label{sec:acervos}

A terceira rodada responde a uma pergunta documental mais exigente: localizar uma referência ao riso ou ao sorriso em cada década de cada um dos treze endereços indicados pelo autor. O protocolo passa a cruzar acervo e década, em vez de permitir que uma ocorrência localizada numa instituição preencha automaticamente a lacuna de outra. A consulta inicial ocorreu em 8 de novembro de 2026. A conferência complementar foi realizada em 9 de novembro. A solicitação ainda não foi integralmente satisfeita: os resultados confirmados e as quotas pendentes são apresentados separadamente no apêndice e nos dados do repositório \repo.

\subsection{O que constitui uma referência e o que constitui uma atestação}

Para esta ampliação, uma referência documental pode ser uma menção lexical a \forma{riso}, \forma{sorriso} ou \forma{gargalhada}, uma rubrica descritiva ou uma grafia mimética. Essas classes respondem a perguntas diferentes. Um poema com a palavra \forma{sorriso} comprova uma referência temática, mas não comprova uma determinada maneira de grafar o som do riso. Apenas rubricas e grafias miméticas entram no núcleo de representações já definido; os demais documentos ficam no inventário complementar. A confirmação exige conferir o trecho na imagem ou no texto primário, além de registrar a procedência e a base de datação. Quando só o ano da pasta foi identificado, isso é declarado, sem atribuir dia ou mês.

Foram recuperados 652 resultados da busca por \forma{risos} na coleção Jornais da BORE/USP e adquiridos 25 PDFs candidatos. Esses resultados incluem formas lexicais relacionadas e não equivalem a 652 ocorrências autênticas de risada. A seleção fornece uma referência conferida para treze décadas: 1870 a 1970, 2000 e 2010. As décadas de 1980 e 1990 permanecem sem referência selecionada nesta busca; não se conclui daí que o acervo esteja vazio ou que tais formas não fossem usadas. Os extremos do catálogo consultado, 1873 e 2012, também não garantem continuidade das séries entre esses anos \cite{r3archiveH04}.

No acervo de Jardinópolis, a seleção abrange as onze décadas do intervalo das pastas consultadas, de 1926 a 2022. Nove referências são lexicais e duas são grafias miméticas. Isso satisfaz a quota temática desse acervo; não equivale a uma grafia onomatopaica em cada década nem à leitura integral de todos os seus jornais. O inventário complementar reúne 74 referências. Há 51 células com confirmação no acervo indicado, 23 com confirmação em fonte alternativa e 61 linhas com quota pendente. Duas linhas dizem respeito a itens individuais com situação documental própria.

\subsection{Novos testemunhos do núcleo}

Na \textit{Imprensa Ytuana}, ano IV, número 183, de 13 de setembro de 1879, a discussão intitulada \textit{Colonisação chineza} reproduz uma intervenção parlamentar seguida de \forma{(Riso)}. A data foi conferida no cabeçalho da página que ocupa a segunda posição do PDF; a primeira posição do arquivo é uma folha sem o texto do jornal. A ocorrência é uma rubrica descritiva e não uma imitação de som \cite{r3doc1}. Sua inclusão documenta uma prática editorial antiga, sem estabelecer qualquer filiação com \forma{kkkk}.

No acervo de Jardinópolis, a página 3 de \textit{A Cidade}, de 9 de janeiro de 2016, traz \forma{Hahaha...} na coluna \textit{Só Rindo}, assinada por Claudio Aguilhera. Na página 5 da edição de 29 de janeiro de 2022, uma entrevista inclui \forma{kkkkk} no interior das respostas. Os cabeçalhos e os trechos foram inspecionados. Trata-se de duas grafias miméticas documentadas na imprensa, em gêneros diferentes: uma coluna de humor e uma entrevista \cite{r3doc18,r3doc19}. A ocorrência de 2022 permite discutir a circulação da convenção digital no jornal, mas não data sua criação e não demonstra uma origem impressa anterior aos bate-papos.

Na edição de 7 de abril de 2025 do \textit{Correio Braziliense}, a entrevista de Walter Lima Jr. realizada por Ricardo Daehn, na seção \textit{Diversão \& Arte}, emprega \forma{(risos)}. O localizador é a página 22 do PDF integral. Esse registro pertence ao jornal de Brasília e não ao \textit{Correio Braziliense} de Hipólito da Costa, iniciado em 1808 \cite{r3doc21,r3archiveH06}. Na terceira rodada, os quatro novos documentos elevaram o núcleo a 23 atestações: quatorze imagens inspecionadas, seis transcrições primárias eletrônicas e três interações reportadas em pesquisas. As referências lexicais complementares não foram somadas a esse total.

\subsection{Limites temporais e institucionais}

Os treze endereços não correspondem a treze séries completas, independentes e atualizadas até 2026. O catálogo do Google News para o \textit{Jornal do Brasil} informa os extremos 8 de abril de 1891 e 31 de dezembro de 1999. O endereço do Senado oferece o \textit{Jornal do Senado} de abril de 1995 a dezembro de 2019 e a revista \textit{Em Discussão!} de abril de 2010 a abril de 2017. Não se preenchem as décadas ausentes desses recortes com documentos de outras coleções do mesmo produtor \cite{r3archiveH07,r3archiveH09}.

Na BBM/USP, a descrição da coleção Gazeta remete à série 1808--1821, mas o item efetivamente recuperado foi um PDF de 28 páginas de 1817. A busca textual nesse arquivo não autenticou uma representação de riso; isso não equivale à análise de todos os números da Gazeta. No APESP, o item D000071 descreve a \textit{Gazeta do Rio de Janeiro}, número 48, de 25 de fevereiro de 1809, com quatro páginas, mas o registro consultado não fornece fac-símile. Informa ainda que o original está fora de circulação por conservação preventiva. A existência bibliográfica fica confirmada; a ocorrência de riso, não \cite{r3archiveH12,r3archiveH13}.

Na Hemeroteca da BN, as faixas do seletor incluem material de Lisboa do século XVIII. Esse extremo não constitui data inicial da imprensa brasileira nem atestação de riso. O leitor do item mais antigo não carregou páginas úteis durante a consulta. Em contrapartida, um PDF público do \textit{Diario Novo}, de 18 de março de 1845, permitiu conferir uma referência lexical aos risos em texto religioso. O documento é preservado como testemunho lexical, sem deslocar automaticamente para 1845 o início do núcleo de grafias miméticas \cite{r3doc20,r3archiveH05}.

A conferência complementar acrescentou seis décadas à seleção da Hemeroteca Digital Brasileira. Foram conferidos \textit{O Malho} (1903), \textit{A Imprensa} (1911), \textit{Carioca} (1950), \textit{Correio da Manhã} (1965), \textit{O Processo} (1978) e \textit{O Liberal} (1989) \cite{acervoD030,acervoD031,anual1950,acervoD032,acervoD033,acervoD034}. Os anos 1920, 1930 e 1940 já estavam representados na matriz. O exemplo de \textit{Carioca} integra também a cobertura anual, como grafia mimética em letra de música publicada no periódico.

No H07, a menção \forma{um largo sorriso} foi localizada no texto indexado de um recorte do \textit{Jornal do Brasil}, de 1977, preservado pelo Arquivo Central da Unicamp. A ocorrência é lexical. Para o H03, a versão eletrônica publicada pela própria \textit{Folha de S.Paulo} em 1994 fornece uma rubrica já documentada no núcleo. A matriz identifica as duas referências como confirmadas em fonte alternativa e mantém registrada a necessidade de conferência no Google News e no Acervo Folha, respectivamente \cite{acervoD036,folha1994}.

A oitava rodada acrescentou cinco décadas à seleção direta da Hemeroteca. O \textit{Araripe} de 1857 registra uma menção a gargalhada. Quatro edições do \textit{Diario de Pernambuco}, de 1865 a 1895, também integram a seleção da Biblioteca Nacional. Essas edições contam para a matriz de H05 e documentam o título associado a H10 em fonte alternativa. Cada exemplar compartilhado tem uma única chave bibliográfica \cite{acervoD050,acervoD038,acervoD040,acervoD042,acervoD044}.

Para o \textit{Diario de Pernambuco}, a busca examinou PDFs da Biblioteca Nacional e publicações do próprio jornal. Foram registradas dez décadas em fontes alternativas. Para a \textit{Folha de S.Paulo} e o \textit{Correio Braziliense}, entrevistas eletrônicas datadas acrescentaram cinco referências. A coluna de resultado conserva a distinção entre a edição eletrônica do jornal e o portal de acervo indicado. As ocorrências selecionadas no texto digitalizado dos PDFs têm página e edição identificadas; a matriz registra essa modalidade de conferência.

Em \textit{O Globo} (H02), foram localizadas rubricas em entrevistas publicadas pelo próprio jornal. A entrevista de Ferreira Gullar com João Cabral de Melo Neto saiu em 27 de setembro de 1987 e foi republicada no Blog do Acervo em 2020. A matriz usa o ano da publicação original informado pelo jornal e identifica a republicação consultada. Os blogs Amplificador e Sonar documentam ocorrências de 2015 e 2020, respectivamente. As três referências foram classificadas como confirmações em fonte alternativa, com a conferência no portal de acervo ainda pendente \cite{acervoD058,acervoD059,acervoD060}.

A busca complementar abrangeu poemas, letras publicadas, folhetins, crônicas, anúncios de cinema, colunas de televisão e reportagens sobre teatro. A data da linha segue a publicação consultada ou a edição original identificada na republicação. O conteúdo sobre uma obra antiga pode documentar a década em que foi publicado. Quando o trecho é uma menção lexical, como a reação de espectadores a uma peça, ele integra o inventário temático.

O leitor indicado pela CEPE permitiu conferir onze décadas diretamente em fac-símiles. A coleção de jornais recifenses do século XIX forneceu quatro referências, de \textit{Academia Popular}, \textit{Academus}, \textit{A Academia} e \textit{O Alpha}. A série do \textit{Diário da Manhã} forneceu sete, dos anos 1920 aos anos 1980, incluindo um anúncio de cinema em 1935 e uma coluna de televisão em 1985 \cite{acervoD068,acervoD069,acervoD070,acervoD071,acervoD062,acervoD067}. Quando a data impressa difere da pasta, a matriz usa a data impressa. A posição da imagem no leitor permanece no localizador.

As edições de 2005 e 2014 do \textit{Jornal do Senado} foram consultadas pelo texto indexado dos PDFs oficiais da Biblioteca Digital do Senado. Os trechos registram a reação do público ao mamulengo e fazem referência à série \textit{Os Aspones}. O Blog do Acervo de \textit{O Globo} também republica uma entrevista de 1991 em que Zilda Cardoso menciona a gargalhada da turma. Essas três linhas identificam a fonte alternativa e a pendência de verificação no portal originalmente indicado \cite{acervoD072,acervoD073,acervoD074}.

Os perfis distinguem datas editoriais, extremos catalográficos, exemplares acessíveis e horizontes ainda a pesquisar. O portal do APESP não teve seus extremos gerais determinados; atribuir-lhe um início arbitrário esconderia uma lacuna de inventário. Os portais de jornais e os mecanismos da CEPE devolveram informações e candidatos, mas continuam com quotas não satisfeitas. Uma busca sem resultado, um OCR deficiente, um registro sem imagem e uma série fora do recorte são situações distintas. Nenhuma delas autoriza fabricar citações.

\subsection{Interpretação estatística da quota}

Seja $h=(a,d)$ uma célula formada pelo acervo $a$ e pela década $d$. A meta documental é $q_h=1$. Definem-se $I_h=1$ quando ao menos uma referência daquela classe foi autenticada no próprio acervo, e $I_h=0$ enquanto a meta não foi satisfeita. Uma eventual razão $\sum_h I_h/H$ descreve cumprimento operacional de quotas, desde que o conjunto de $H$ células elegíveis esteja previamente delimitado. Não estima a frequência histórica do riso, a popularidade de uma grafia ou sua probabilidade de origem. Os recortes ainda não inventariados impedem interpretar o total de linhas da matriz como universo documental completo.

A quota unitária foi escolhida para atender ao desenho documental solicitado pelo autor. Não foi obtida por um cálculo de precisão amostral. Com apenas um documento por célula, seleção dirigida e probabilidades de inclusão desconhecidas, não há base para intervalos de confiança de prevalência. Para investigar frequência seria necessário definir unidades, denominadores, procedimento de sorteio e esforço de inspeção, conforme o planejamento discutido na metodologia. Também não se contam como observações independentes duas cópias do mesmo número de jornal mantidas em instituições diferentes. A matriz conserva a procedência de acesso; a análise histórica deve identificar e tratar duplicatas.

O repositório \repo\ disponibiliza \nolinkurl{archive_profiles.csv}, \nolinkurl{archive_decade_references.csv}, \nolinkurl{archive_decade_matrix.csv} e \nolinkurl{archive_summary.json}, além do gerador \nolinkurl{build_method_inventory.py}. Os arquivos tornam verificáveis tanto as referências confirmadas quanto o caráter ainda incompleto da cobertura por acervo. As imagens e os PDFs adquiridos permanecem locais quando não há licença de redistribuição verificada; URLs, localizadores e resumos de integridade permitem reencontrar as fontes.

\section{Seleção cronológica de obras, periódicos e páginas inéditos}

A quarta rodada amplia o levantamento a partir de três endereços indicados pelo autor: a coleção geral da Brasiliana/USP, uma captura da página de entrada do arquivo de comunidades do Orkut e o portal da Wayback Machine. Adotou-se uma seleção cronológica de títulos ainda não examinados, incluindo sites, fóruns e blogs, conforme delimitação expressa do autor. A unidade editorial é uma referência por obra, periódico ou site. Essa regra difere da quota por acervo e década da rodada anterior, cujo estado de conclusão permanece publicado. Não se reapresentam títulos anteriores como novidades por terem sido reencontrados em outro repositório.

\subsection{Seleção, exclusão de repetições e cronologia}

O inventário anterior de obras e jornais, incluindo as indicações auditadas sem ocorrência confirmada, foi conservado como lista de exclusão. Os títulos selecionados foram confrontados com esse inventário; volumes de uma mesma obra e números de um periódico não foram multiplicados para aumentar o total. A busca dirigida por \forma{riso} na Brasiliana retornou 1.078 itens na consulta de descoberta registrada. Resultados ordenados por data foram usados para descoberta, com inspeção de exemplares escolhidos; não se leu integralmente esse conjunto. O mesmo identificador de coleção fornecido pelo autor foi acessado no endereço atual \url{https://digital.bbm.usp.br/handle/bbm/1}. A falha de certificado no endereço antigo não foi resolvida mediante desativação da verificação de segurança.

A seleção executada reúne dezessete títulos: quatorze da Brasiliana, um blog, um fórum e uma comunidade pública do Orkut. O conjunto não constitui um contínuo anual: o intervalo dos testemunhos selecionados vai de 1645 a 2017, e a data da consulta é 2026. O exemplar de 1645 é o mais antigo confirmado nesta seleção, não a primeira ocorrência portuguesa de \forma{riso} nem o livro mais antigo da biblioteca. A existência de materiais quinhentistas no catálogo não permite atribuir-lhes a ocorrência sem recuperar o trecho.

As referências complementares distribuem-se em treze menções lexicais, dois registros lexicográficos e duas grafias miméticas. As duas últimas integram o núcleo de representações; sua inclusão antecedeu a ampliação anual. No balanço consolidado, o núcleo contém 121 registros: 23 E1, 95 E2 e três E3. As demais referências documentam usos, sentidos ou tematizações do riso; não são convertidas em exemplos de \forma{kkkk}.

\subsection{Séculos XVII e XVIII: leitura das fontes e história da escrita}

Na \textit{Chronica da Companhia de Iesv na provincia de Portvgal}, volume I, de Baltasar Teles, edição de 1645, a página 266 do PDF contém uma referência a risos e zombarias. A leitura foi conferida na imagem. O OCR também oferece um falso candidato em outra página, lendo como \forma{risos} uma sequência que não autentica o fenômeno. No trecho selecionado, a tipografia emprega o chamado \emph{s longo}, reconhecido mecanicamente como \forma{f}. A normalização para \forma{risos} é uma decisão de transcrição, não uma alegação de mudança fonética de \forma{rifos} para \forma{risos} \cite{r4chronicacompanhiaiesvportugal}.

O volume VII do \textit{Vocabulario Portuguez e latino}, de Rafael Bluteau, de 1720, apresenta o verbete \forma{Riso} na página impressa 345, correspondente à página 353 do PDF. O verbete aborda sinais corporais e manifestações da alegria, mas suas explicações devem ser lidas como história da lexicografia e das ideias de sua época, sem reproduzir sua fisiologia como ciência atual. No \textit{Diccionario da lingua portugueza} reformado por Antônio de Moraes Silva, de 1789, o volume II contém o verbete \forma{RISO}, na página 354 do PDF. São obras lexicográficas relacionadas, não duas provas independentes de uma mesma genealogia digital \cite{r4vocabularioportuguezlatino,r4diccionariolinguaportuguezamorais}.

Em \textit{Viola de Lereno}, volume I, de 1798, a expressão \forma{Mostra em riso moderado} ocorre na página 259 do PDF. A passagem descreve o riso e a exposição dos dentes; não grafa uma risada por repetição de sílabas. A seleção mantém apenas uma referência para a obra, em vez de somar seus volumes e reedições como títulos novos \cite{r4violadelereno}.

\subsection{Periódicos, teatro, poesia e prosa dos séculos XIX e XX}

O recorte inclui três periódicos inéditos no inventário anterior: \textit{O Patriota}, no conjunto de números de 1813; \textit{Gazeta Pernambucana}, de 1822; e \textit{Nitheroy}, de 1836. As expressões selecionadas, respectivamente, remetem a ser digno de riso, a merecer risada e a um riso de alegria. São testemunhos lexicais localizados nas páginas 25, 26 e 140 dos PDFs correspondentes. A \textit{Gazeta Pernambucana} não foi confundida com o \textit{Diario de Pernambuco} nem com a \textit{Gazeta do Rio de Janeiro} \cite{r4opatriota,r4gazetapernambucana,r4nitheroy}.

Nas \textit{Reflexões dramaticas para uso dos canditados que se dedicam a scena}, de João Caetano, edição de 1837, a recomendação de evitar rir às gargalhadas em determinado papel aparece na página 22 do PDF, página impressa 16. Trata-se de reflexão sobre representação cênica, não de proibição gramatical geral. \textit{Primeiros cantos}, de Gonçalves Dias, edição de 1846, oferece uma referência lexical na página impressa 47, PDF 51. \textit{Gennesco}, de 1866, \textit{A fallencia}, de 1902, e \textit{Negrinha}, de 1920, acrescentam usos em prosa com referências a risos, gargalhadas e reações de personagens. As imagens dos trechos foram inspecionadas, sem considerar o texto de OCR uma transcrição automaticamente fidedigna \cite{r4reflexoesdramaticas,r4primeiroscantos,r4gennesco,r4afallencia,r4negrinha}.

Em \textit{Poesia, 1958/1978}, de Ana Hatherly, edição de 1980, a página 35 do PDF contém \forma{O riso tem lágrimas}. O intervalo do título não permite atribuir esse poema a um ano específico entre 1958 e 1978 sem outro cotejo. O registro foi datado pelo exemplar consultado. Em \textit{Rubens Borba e Moraes: Anotações de um Bibliófilo}, de Cristina Antunes, publicação de 2017, a página impressa 144, PDF 146, reproduz uma referência a poesias que serviam de riso. A redação remete a material antigo, mas o testemunho consultado é a publicação de 2017; não se retrodata automaticamente a frase a um manuscrito anterior \cite{r4poesia19581978,r4rubensborbamoraesanotacoesbibliofilo}.

\subsection{Wayback, blog e fórum: três datas que não devem ser confundidas}

Uma página preservada pode oferecer três datas diferentes: a data que o conteúdo exibe, a data de criação do espaço que o hospeda e a data da captura. Além disso, solicitar um instante à Wayback pode levar a outro registro disponível. O manifesto conserva tanto a URL solicitada quanto a URL efetivamente servida. Essa diferença foi decisiva na triagem de uma página incorporada do \textit{Humor Tadela}: uma consulta iniciada numa captura de dezembro de 1998 conduziu a um recurso de novembro de 2002, sem ocorrência pertinente autenticada. O recurso não foi contado como risada documentada em 1998.

No \textit{Kibe Loco}, a postagem \textit{Uma chance (parte 2)} exibe a data de 27 de junho de 2007 e a menção \forma{não contendo o riso}. O conteúdo foi recuperado numa captura efetiva de 1º de julho de 2007, embora o pedido inicial ao endereço de entrada estivesse associado a 10 de junho. A referência comprova a presença lexical naquela postagem e não a publicação em junho de conteúdo que a própria página data posteriormente \cite{r4kibeloco}.

No fórum Cifra Club, a captura de 15 de setembro de 2010 do tópico \textit{como afinar meio tom abaixo} preserva uma resposta com \forma{kkkkkkkkkkkkk}, identificada na interface como enviada em 12 de julho de 2010, às 15h21. O tópico foi iniciado em novembro de 2003. A data de abertura não pode ser transferida à resposta acrescentada em 2010. O registro fornece texto primário eletrônico preservado de escrita com \forma{k} repetido; não antecipa para 2003 a ocorrência mais antiga deste levantamento \cite{r4forumcifraclub}.

\subsection{Orkut: página de entrada, comunidade e descrição preservada}

A captura indicada pelo autor, de 1º de outubro de 2014, apresenta a entrada do arquivo de comunidades do Orkut. A página descreve o período de funcionamento do serviço, de janeiro de 2004 a setembro de 2014. Esse texto institucional não é uma conversa de 2004 e não comprova sozinho uma forma específica de riso \cite{r4orkutindex}.

A navegação pelo índice levou à comunidade pública \textit{Kiko Dança Psy-TRANCE}. Sua descrição contém \forma{hehehe} na captura efetiva de 4 de outubro de 2014. Embora a página indique criação da comunidade em 24 de fevereiro de 2006, a descrição pode ter sido editada posteriormente. Na ausência de uma versão anterior do trecho, o testemunho foi datado pela captura de 2014, estabelecendo apenas que a sequência estava preservada até esse momento. Não se atribui a redação a 2006, nem se apresenta o exemplo como prova de uma primeira utilização no Orkut \cite{r4orkutcomunidades}.

Os três espaços digitais cumprem papéis diferentes nesta seleção: referência lexical em blog, resposta de fórum com marcador de riso e descrição de comunidade com grafia mimética. Sua comparação esclarece funções e condições de preservação, mas não reconstrói uma cadeia causal entre plataformas. Uma captura arquivada tampouco representa todas as interações que o serviço abrigou.

\subsection{Reprodução e limites da ampliação}

O repositório \repo\ publica a lista de exclusão anterior, a tabela de títulos inéditos e o resumo da rodada em \nolinkurl{revision4_exclusions.csv}, \nolinkurl{revision4_new_titles.csv} e \nolinkurl{revision4_summary.json}. As aquisições têm localizadores e SHA-256 no manifesto geral. A lista cronológica aparece no apêndice. A seleção não é um censo da Brasiliana ou da web preservada, não é equivalente à seleção por quota de acervo e década e não permite estimar prevalências. Não se aplicou uma fórmula amostral para justificar dezessete títulos: o número é o resultado desta seleção dirigida, condicionada pela novidade em relação ao inventário, pela diversidade temporal e pela autenticação dos trechos.

\section{Discussão: do impresso às redes digitais}
\subsection{O que se pode afirmar sobre a origem de \forma{kkkk}}
As fontes inspecionadas sustentam uma história de representações múltiplas do riso. Não sustentam uma linha necessária \forma{quiá} $>$ \forma{cá} $>$ \forma{quá} $>$ \forma{kakaka} $>$ \forma{kkkk}. Para testar essa linha, seriam úteis documentos intermediários, declarações contemporâneas de adaptação, padrões de circulação, reedições, tradições de quadrinhos e corpora digitais datados. Uma série de setas não substitui esses elos.

Também não se demonstra independência absoluta. Uma geração não precisa ter lido conscientemente \textit{Til} para participar de uma convenção transmitida por quadrinhos, escola, família ou mídia. Transmissão indireta é uma possibilidade; reinvenção também. O exemplo hipotético do material inicial, em que leitores adotariam uma forma dos livros e a passariam adiante, é uma hipótese legítima. Seu requisito não deve ser estreitado a fala cotidiana nem a manutenção de um radical identificável.

A formulação sustentada é: há antecedentes impressos funcional e parcialmente formalmente semelhantes ao riso digital; a coleta não estabelece uma cadeia específica de descendência, tampouco sua inexistência. A forma digital está documentada no conjunto acadêmico ligado a 2007.2 e pode ser anterior. O intervalo ``1998 a 2002'' não foi validado como origem. Atribuí-lo a consenso de historiadores ou linguistas, sem estudos identificados, seria um argumento de autoridade sem referência.

\subsection{IRC, mIRC, BBS, salas públicas e mensageiros}
O material enumera mIRC, ICQ, bate-papos de UOL, BOL e Terra, MSN, Orkut e jogos. Esses ambientes são relevantes para investigar a difusão de convenções digitais, mas precisam ser diferenciados. IRC é um protocolo ou ambiente de conversação; mIRC é um cliente. Mensageiros privados, fóruns, serviços de recados e salas públicas têm formas distintas de arquivamento, público e temporalidade. ``Forjado sequencialmente em três plataformas'' exige evidência de transferência entre elas.

O levantamento não recuperou logs autenticados de 1995--1997 que sustentem a alegação de ausência, nem uma coleção de 1999--2000 que demonstre \forma{kakaka} como estágio obrigatório. A existência histórica dessas plataformas não comprova a presença de uma grafia específica. Um testemunho retrospectivo pode orientar buscas, mas está sujeito a memória reconstrutiva. Sua data é a do depoimento, salvo quando acompanhado de arquivo contemporâneo verificável.

As alegações de massificação pelo MSN e de uso em jogos como \textit{Ultima Online}, \textit{Tibia} e \textit{Counter-Strike} permanecem hipóteses de difusão. A afirmação de banimentos frequentes por confusão com uma sigla supremacista precisaria de regras de moderação, mensagens administrativas ou registros datados. Sem esses documentos, não se transforma uma anedota plausível em prova de origem.

\subsection{A hipótese ergonômica e seus limites causais}
Repetir uma tecla pode ser mecanicamente mais simples que alternar duas. O nome \forma{cá} também oferece uma interpretação possível da cadeia de letras. Isso torna plausível uma hipótese de economia gráfica, mas não prova a motivação dos primeiros usuários. A escrita digital contém inúmeras formas mais longas e custosas, como \forma{huahuahua} e sequências heterogêneas, empregadas por seus efeitos expressivos. A redução de esforço é um fator possível entre outros.

A hipótese poderia ser testada experimentalmente com tempos de digitação, familiaridade, dispositivo e tarefa controlados, ou documentalmente com entrevistas contemporâneas e séries de uso. Uma experiência atual demonstraria diferenças atuais de custo, não a causa histórica original. Também se deveria distinguir repetição por toque, pressionamento contínuo e recursos de autocompletar: diferentes sistemas e teclados comportam-se de maneiras distintas.

A presença ou ausência de \forma{k} no inventário alfabético normativo não determina sua disponibilidade em máquinas ou na escrita. Letras podem circular em nomes, símbolos e empréstimos antes de mudanças normativas. A reforma ortográfica não é uma evidência de datação de \forma{kkkk}. O argumento de que ``a letra só entrou no alfabeto depois'' confunde norma ortográfica, práticas sociais e existência material do caractere.

\subsection{Tecnologia, imprensa e coexistência de convenções}
Tipografia, composição editorial, espaço de coluna e práticas de tradução podem influenciar o texto. Entretanto, não se deve atribuir a cada década uma personalidade coletiva: riso ``aristocrático'' em 1900, ``malandro'' em 1960 e ``funcional'' em 1980 são interpretações que precisariam de indicadores. Uma passagem de folhetim, uma rubrica parlamentar e um balão de quadrinhos pertencem a gêneros distintos, mesmo quando publicados no mesmo ano.

A explicação de que a censura dos anos 1940 ou 1970 substituiu onomatopeias por rubricas exigiria comparar publicações, editorias e períodos sob procedimentos definidos. Correlação temporal com um regime político não estabelece causalidade gráfica. Do mesmo modo, a popularidade da Disney não demonstra que o \forma{quá} brasileiro se originou ali; a capa de 1935 fornece um testemunho anterior ao marco proposto. A força de uma indústria cultural é compatível com difusão de formas já existentes.

As hipóteses sobre origem de \forma{rs} em telegramas, gabinetes e redações merecem investigação, não endosso automático. Seriam necessários manuais, cartas, rascunhos ou laudas datadas com interpretação inequívoca. A facilidade de explicar uma abreviação a partir de \forma{risos} não fornece sua data inicial. Até que se recupere esse material, a genealogia específica ``máquina de escrever $>$ chat'' permanece indeterminada.

\section{Variação pragmática contemporânea e o problema das gerações}
\subsection{Maiúsculas, extensão e sequências heterogêneas}
\forma{kkkk}, \forma{KKKKKK}, \forma{hehe}, \forma{hahaha}, \forma{rsrs} e \emph{keysmashes} podem diferenciar intensidade, afiliação, ironia, constrangimento ou estilo. A interpretação depende de posição na conversa, relação entre interlocutores, gênero discursivo, plataforma e repertório local. Maiúsculas e alongamento podem intensificar a expressão, mas não há fundamento para uma regra universal segundo a qual minúsculas indiquem falsidade e maiúsculas sinceridade.

Uma cadeia heterogênea como \forma{asdfghjkl} pode figurar descontrole expressivo e funcionar como prática aprendida. Não precisa ter surgido com telas de toque: sequências baseadas no alinhamento das teclas são compatíveis com teclado físico. Um formato pode ser empregado em diferentes dispositivos, e uma explicação causal precisa datar sua ocorrência. Não se infere o aparelho utilizado apenas pela forma textual.

\subsection{Emoticons, emojis e respostas verbais}
Emoticons como \texttt{:-)} e emojis são recursos distintos. Os primeiros combinam caracteres de escrita; os segundos são sinais pictográficos codificados, com desenho que pode variar entre plataformas. Emoji de caveira, rosto chorando e rosto com lágrimas de alegria podem adquirir leituras humorísticas, mas o valor convencional não se esgota numa tradução fixa. O uso de caveira para humor pode explorar a expressão ``morrer de rir''; isso não prova que toda ocorrência de caveira denote gargalhada.

O repertório também inclui respostas como \forma{rindo}, \forma{rindo de nervoso}, comentários lacônicos, imagens e figurinos visuais. Uma resposta com ponto final pode comunicar secura ou ironia em determinado grupo, sem que o ponto carregue sempre esse significado. A revisão de Badrakh e colaboradores \cite{badrakh2024} destaca a relevância e a ambiguidade dos sinais digitais em brincadeiras online. A interpretação precisa da sequência conversacional; uma lista de equivalências descontextualizadas perde esse componente.

O significado de um marcador pode mudar sem alteração em sua forma. Esse é um caso pertinente de mudança semântica ou pragmática. Portanto, o fato de não haver derivação de um radical não exclui mudança linguística. A expansão de leituras humorísticas de um signo visual permite precisamente investigar convencionalização, inferência e circulação histórica.

\subsection{Geração Z e Alpha: hipóteses, não gramáticas uniformes}
As respostas iniciais atribuem à geração Z e à Alpha preferências uniformes, com limites de nascimento fixos e oposição rígida aos millennials. O presente corpus não contém amostra brasileira por idade que permita estimar tais preferências, e não inclui dados próprios sobre crianças da geração Alpha. Identificar uma tendência num grupo online não autoriza generalização a todos os integrantes de uma coorte.

Kempe e Raviv \cite{kempe2025} estudam a convencionalização de emojis faciais e não encontram as diferenças geracionais previstas naquele desenho. O resultado não elimina toda variação etária, cultural ou contextual, mas impede tratar manchetes sobre ``emoji de velho'' como regra científica universal. A revisão mais ampla de Costanza, Rudolph e Zacher \cite{costanza2023} também alerta para a fragilidade de atribuições geracionais que confundem idade, período e coorte.

Uma investigação brasileira precisaria recrutar participantes com idades conhecidas, caracterizar região e comunidades de prática, medir usos e interpretações em tarefas comparáveis e analisar diferenças de plataforma. Para menores de idade, seriam necessários procedimentos éticos apropriados ao desenho. O corpus documental deste artigo não equivale a uma pesquisa com participantes, e não se declara aprovação ética inexistente.

\section{Bibliografia, autoridade e integridade da evidência}
\subsection{Referências da UnB e atualização teórica}
A escolha de bibliografia não se faz apenas pela data mais recente. Obras de referência podem continuar relevantes; artigos recentes são indispensáveis para debates e dados novos. O Edital 04/2026 da UnB fornece uma base institucional verificável para Martelotta e Seara, Nunes e Lazzarotto-Volcão \cite{unb2026}. A bibliografia inclui ainda estudos recentes sobre discurso e multimodalidade. Não se presume que a lista seletiva esgote o ensino ou a pesquisa do departamento.

Campbell \cite{campbell2021} oferece enquadramento histórico-comparativo contemporâneo; Ilari \cite{ilari2018}, uma referência para a formação românica. Monte \cite{monte2019} é especialmente importante porque associa discussão de tratamento a documentação datada. Iida e Akita \cite{iida2024} atualizam a relação entre experiência e iconicidade. Os estudos sobre sinais digitais e emojis ampliam a interpretação pragmática sem fornecer, indevidamente, a origem do \forma{kkkk} brasileiro.

\subsection{Revisão das referências sugeridas no material inicial}
Saussure, Coseriu, Ullmann, Bréal, Jakobson, Tarallo, Faria, Christiansen e Chater são nomes associados a tradições relevantes. Contudo, sua autoridade não autentica uma tese específica sem passagem pertinente. Não se atribui a Guimarães a demonstração de que onomatopeias não possuem raízes históricas, nem a Jakobson uma regra de independência de todas as imitações sonoras. Títulos, edições e usos argumentativos precisam ser conferidos antes da inclusão como sustentação.

O material atribui a Ilari um livro denominado \textit{Linguística histórica}, com dados bibliográficos não corroborados na coleta; a obra identificada é \textit{Linguística românica}. Esse tipo de correção não é cosmético: permite que outro leitor encontre a fonte. Também é inadequado transformar um livro sobre fonética do latim em autoridade suficiente para a história de qualquer palavra portuguesa sem localizar a discussão correspondente.

Uma referência acadêmica recente pode repetir uma hipótese de divulgação sem fornecer nova prova primária. Identificou-se, por exemplo, um artigo de 2025 que remete a uma página de divulgação ao narrar a relação entre \forma{quiá} e \forma{kkkk}. Esse encadeamento mostra circulação da narrativa, não confirmação independente da genealogia. Contar citações dessa versão como testemunhos distintos produziria uma falsa multiplicação de evidências.

\subsection{Fidelidade ao material de partida e transparência de autoria}
Ser fiel às questões enviadas exige preservar o seu problema, suas alternativas e seus exemplos, mas não manter como verdade o que a verificação não sustenta. Por isso, o artigo conserva a discussão sobre raiz, mudança, onomatopeia, suporte, periódicos, plataformas e gerações; reclassifica as supostas provas; corrige conceitos e registra lacunas. As citações sem autenticação aparecem como alvos de auditoria, e não como falas efetivas dos autores históricos.

Murillo Edson de Carvalho Souza é o único autor identificado. A elaboração deste manuscrito e dos scripts utilizou assistência de inteligência artificial para organização, redação, programação e busca documental. Ferramentas automatizadas não são coautoras nem fontes históricas. A publicação do material não equivale a revisão humana independente ou avaliação por pares; a responsabilidade pela versão científica submetida e por suas conclusões permanece com o autor. O artigo não herda do template as declarações, afiliações, dados ou coautoria do texto que nele vinha exemplificado.

\section{Limitações e agenda de pesquisa}
O corpus é pequeno, dirigido por problemas e desigual no tempo. As lacunas de várias décadas do século XX impedem estabelecer predomínio de formas, taxas de substituição ou uma cronologia nacional de difusão. A coleta não autenticou a série de 22 transcrições enviada inicialmente e não identificou a primeira ocorrência absoluta de \forma{kkkk}. Esses limites são resultados da auditoria, não espaços a preencher com frases plausíveis.

O acesso digital introduz filtros: preservação seletiva, qualidade de digitalização, OCR, disponibilidade de catálogo, direitos e barreiras técnicas. Uma ausência num acervo pode refletir qualquer desses fatores. A data de uma transcrição eletrônica não é a data da criação literária; a data de publicação de uma dissertação não é necessariamente a data de uma mensagem. O estudo registra essas diferenças, mas não resolve todos os problemas de transmissão editorial dos textos consultados.

A agenda de ampliação tem quatro frentes. Primeiro, construir uma moldura anual de edições e números de periódicos, com países, gêneros e extensões. Segundo, buscar registros digitais contemporâneos anteriores a 2007, validando procedência e data, sem confundir memórias recentes com logs antigos. Terceiro, estudar canais intermediários como teatro, quadrinhos, cartas e livros escolares, capazes de testar continuidade indireta. Quarto, investigar a pragmática contemporânea com dados brasileiros e variáveis sociais observadas, incluindo comparações entre dispositivos e comunidades de prática.

A aplicação do plano estatístico dependerá do objetivo: existência, detecção de formas raras ou estimação de frequência. As quotas de três livros e três periódicos podem permanecer como seleção de exemplos para leitura; não devem ser apresentadas como amostra representativa. Um estudo probabilístico requer sorteio, denominadores e análise compatíveis. Um estudo filológico pode ter grande valor com poucos testemunhos, desde que não extrapole sua abrangência.

\section{Conclusão}
Os testemunhos recuperados confirmam diversidade de formas de escrever o riso, incluindo \forma{quiá}, \forma{cá}, \forma{ah} e \forma{quá}, além de cadeias digitais de \forma{k}. O registro de \textit{O Tico-Tico}, de 1935, fornece uma localização verificável de \forma{quá}; a dissertação de Young oferece documentação de \forma{kkkk} em pesquisa situada em 2007.2. Esses achados sustentam atestações específicas, sem determinar a origem absoluta da convenção digital.

A contribuição conceitual é rejeitar duas inferências simétricas. A semelhança entre representações não demonstra uma genealogia; o caráter onomatopaico e as lacunas documentais não demonstram independência. Onomatopeias podem integrar a história das línguas, e mudanças gráficas podem ter consequências linguísticas. A questão é empírica: quais canais, documentos e regularidades sustentam uma relação proposta?

O dimensionamento estatístico mostra por que poucas referências por período servem melhor à exemplificação que à prova de ausência ou predomínio. A documentação pública, as tabelas de cobertura e os scripts permitem ampliar e corrigir a pesquisa. A interpretação cientificamente defensável permanece aberta à descoberta de novos registros, inclusive anteriores aos aqui localizados.

\section*{Disponibilidade de dados e materiais}
Manuscrito, bibliografia, inventários, códigos de evidência e scripts: \repo. Os downloads permitidos para distribuição acompanham o repositório; arquivos com condições de republicação não verificadas permanecem no acervo local de trabalho, com links públicos de aquisição. O arquivo \texttt{README.md} descreve a reprodução e as lacunas. A versão consultada deve ser identificada por commit ou release.

\section*{Declarações}
\textbf{Autoria:} Murillo Edson de Carvalho Souza, único autor. \textbf{Financiamento e conflitos de interesse:} não foram informados no material fornecido; nenhuma declaração de inexistência é presumida. \textbf{Assistência automatizada:} utilizada na elaboração e documentação, conforme descrito na discussão. \textbf{Situação editorial:} manuscrito de pesquisa exploratória, sem indicação de avaliação por pares concluída.

\bibliographystyle{plain}\bibliography{references,annual}
\clearpage\appendix\section{Auditoria das supostas citações literais}
A tabela preserva a identificação e a grafia atribuídas às fontes no material inicial. Não reproduz as frases como citações autênticas. O diagnóstico de não autenticação indica insuficiência de comprovação nesta coleta, não demonstração de inexistência em qualquer edição. A matriz CSV conserva uma âncora curta para buscas futuras.
\small\begin{longtable}{p{0.09\textwidth}p{0.27\textwidth}p{0.19\textwidth}p{0.29\textwidth}}
\caption{Os 22 alvos da série enviada: grafias alegadas, sem autenticação das frases.}\\
\toprule Ano & Fonte alegada & Grafia alegada & Falta documental \\\midrule\endfirsthead
\toprule Ano & Fonte alegada & Grafia alegada & Falta documental \\\midrule\endhead
1902 & Canaã & Ah! Ah! Ah! & Sem edição e página. \\
1904 & Gazeta de Notícias, 25/07 & Ah! Ah! Ah! & Sem página e autoria da crônica. \\
1911 & Triste fim de Policarpo Quaresma & Ah! ah! & Sem edição e página; separar folhetim de 1911 e livro de 1915. \\
1912 & O Malho, 12/10 & Quá! Quá! Quá! & Sem página e número. \\
1928 & Macunaíma & Ah! ah! ah! & Sem página e edição. \\
1925 & Correio da Manhã, 14/02 & Ah! ah! ah! & Sem página e autoria. \\
1934 & São Bernardo & Ah, ah! & Sem página e edição. \\
1935 & Fon-Fon!, 08/06 & Quá! quá! quá! & Sem página. \\
1946 & Sagarana & He! he! he! & Sem conto, página e edição. \\
1944 & Diário de Notícias, 18/03 & (Risos na assistência) & Sem página e identificação do ministro. \\
1953 & Almanaque do Pato Donald & Quá-quá-quá! & Sem número, história e página. \\
1955 & Tribuna da Imprensa & He-he-he & Sem dia, página e título da coluna. \\
1967 & Quarup & Ha, ha, ha & Sem página e edição. \\
1969 & O Pasquim, n. 12 & (hahaha) & Sem página e entrevista. \\
1973 & As meninas & Ha-ha-ha! & Sem página e edição. \\
1975 & Revista do Cebolinha & Hi-hi-hi! & Sem número, história e página. \\
1982 & Feliz ano velho & Ha ha ha & Sem página e edição. \\
1986 & Folha de S.Paulo, Painel, 14/11 & Rs. & Sem página. \\
1993 & Veja, Páginas Amarelas, 07/05 & (Risos) & Entrevistado não identificado. \\
1997 & Folhateen, junho & rsrsrs / hahaha & Sem dia, página e reportagem. \\
2005 & Os segredos do Orkut & kakakaka & Sem autor, editora e edição identificados. \\
2007 & Capricho, setembro & KKKKKKKKKK & Sem número, página e reportagem. \\
\bottomrule\end{longtable}\normalsize

\section{Cobertura anual e distinção entre lacuna e ausência}
Os anos abaixo pertencem ao núcleo 1900--2009. A coluna de atestações conta documentos incluídos no inventário, não todos os arquivos baixados nem o universo publicado. O símbolo -- significa que não há atestação confirmada naquele ano nesta coleta; não representa frequência histórica zero. As extensões de 1872 a 2026 constam no arquivo público.
\small\begin{longtable}{rrp{0.65\textwidth}}
\caption{Cobertura anual do núcleo 1900--2009. Não é série de prevalência.}\\
\toprule Ano & Documentos & Interpretação \\\midrule\endfirsthead
\toprule Ano & Documentos & Interpretação \\\midrule\endhead
1900 & 1 & Atestação confirmada \cite{anual1900}. \\
1901 & 1 & Atestação confirmada \cite{queiros1901}. \\
1902 & 1 & Atestação confirmada \cite{anual1902}. \\
1903 & 1 & Atestação confirmada \cite{anual1903}. \\
1904 & 1 & Atestação confirmada \cite{anual1904}. \\
1905 & 1 & Atestação confirmada \cite{anual1905}. \\
1906 & 1 & Atestação confirmada \cite{anual1906}. \\
1907 & 1 & Atestação confirmada \cite{anual1907}. \\
1908 & 1 & Atestação confirmada \cite{anual1908}. \\
1909 & 1 & Atestação confirmada \cite{anual1909}. \\
1910 & 1 & Atestação confirmada \cite{anual1910}. \\
1911 & 1 & Atestação confirmada \cite{anual1911}. \\
1912 & 1 & Atestação confirmada \cite{anual1912}. \\
1913 & 1 & Atestação confirmada \cite{anual1913}. \\
1914 & 1 & Atestação confirmada \cite{anual1914}. \\
1915 & 1 & Atestação confirmada \cite{anual1915}. \\
1916 & 1 & Atestação confirmada \cite{anual1916}. \\
1917 & 1 & Atestação confirmada \cite{almada1917}. \\
1918 & 1 & Atestação confirmada \cite{anual1918}. \\
1919 & 1 & Atestação confirmada \cite{lobato1919}. \\
1920 & 1 & Atestação confirmada \cite{anual1920}. \\
1921 & 1 & Atestação confirmada \cite{anual1921}. \\
1922 & 1 & Atestação confirmada \cite{anual1922}. \\
1923 & 1 & Atestação confirmada \cite{anual1923}. \\
1924 & 1 & Atestação confirmada \cite{anual1924}. \\
1925 & 1 & Atestação confirmada \cite{portugal1925}. \\
1926 & 1 & Atestação confirmada \cite{anual1926}. \\
1927 & 1 & Atestação confirmada \cite{anual1927}. \\
1928 & 1 & Atestação confirmada \cite{diarionacional1928}. \\
1929 & 1 & Atestação confirmada \cite{anual1929}. \\
1930 & 1 & Atestação confirmada \cite{anual1930}. \\
1931 & 1 & Atestação confirmada \cite{anual1931}. \\
1932 & 1 & Atestação confirmada \cite{anual1932}. \\
1933 & 1 & Atestação confirmada \cite{anual1933}. \\
1934 & 1 & Atestação confirmada \cite{anual1934}. \\
1935 & 1 & Atestação confirmada \cite{ticotico1935}. \\
1936 & 1 & Atestação confirmada \cite{anual1936}. \\
1937 & 1 & Atestação confirmada \cite{anual1937}. \\
1938 & 1 & Atestação confirmada \cite{anual1938}. \\
1939 & 1 & Atestação confirmada \cite{anual1939}. \\
1940 & 1 & Atestação confirmada \cite{anual1940}. \\
1941 & 1 & Atestação confirmada \cite{carioca1941}. \\
1942 & 1 & Atestação confirmada \cite{anual1942}. \\
1943 & 1 & Atestação confirmada \cite{anual1943}. \\
1944 & 1 & Atestação confirmada \cite{anual1944}. \\
1945 & 1 & Atestação confirmada \cite{anual1945}. \\
1946 & 1 & Atestação confirmada \cite{anual1946}. \\
1947 & 1 & Atestação confirmada \cite{jornalnoticias1947}. \\
1948 & 1 & Atestação confirmada \cite{anual1948}. \\
1949 & 1 & Atestação confirmada \cite{anual1949}. \\
1950 & 1 & Atestação confirmada \cite{anual1950}. \\
1951 & 1 & Atestação confirmada \cite{anual1951}. \\
1952 & 1 & Atestação confirmada \cite{anual1952}. \\
1953 & 1 & Atestação confirmada \cite{anual1953}. \\
1954 & 1 & Atestação confirmada \cite{anual1954}. \\
1955 & 1 & Atestação confirmada \cite{camara1955}. \\
1956 & 1 & Atestação confirmada \cite{anual1956}. \\
1957 & 1 & Atestação confirmada \cite{anual1957}. \\
1958 & 1 & Atestação confirmada \cite{anual1958}. \\
1959 & 1 & Atestação confirmada \cite{anual1959}. \\
1960 & 1 & Atestação confirmada \cite{anual1960}. \\
1961 & 1 & Atestação confirmada \cite{anual1961}. \\
1962 & 1 & Atestação confirmada \cite{anual1962}. \\
1963 & 1 & Atestação confirmada \cite{anual1963}. \\
1964 & 1 & Atestação confirmada \cite{anual1964}. \\
1965 & 1 & Atestação confirmada \cite{senado1965}. \\
1966 & 1 & Atestação confirmada \cite{anual1966}. \\
1967 & 1 & Atestação confirmada \cite{anual1967}. \\
1968 & 1 & Atestação confirmada \cite{anual1968}. \\
1969 & 1 & Atestação confirmada \cite{anual1969}. \\
1970 & 1 & Atestação confirmada \cite{anual1970}. \\
1971 & 1 & Atestação confirmada \cite{anual1971}. \\
1972 & 1 & Atestação confirmada \cite{anual1972}. \\
1973 & 1 & Atestação confirmada \cite{anual1973}. \\
1974 & 1 & Atestação confirmada \cite{anual1974}. \\
1975 & 1 & Atestação confirmada \cite{anual1975}. \\
1976 & 1 & Atestação confirmada \cite{anual1976}. \\
1977 & 1 & Atestação confirmada \cite{anual1977}. \\
1978 & 1 & Atestação confirmada \cite{anual1978}. \\
1979 & 1 & Atestação confirmada \cite{correio1979}. \\
1980 & 1 & Atestação confirmada \cite{anual1980}. \\
1981 & 1 & Atestação confirmada \cite{anual1981}. \\
1982 & 1 & Atestação confirmada \cite{anual1982}. \\
1983 & 1 & Atestação confirmada \cite{anual1983}. \\
1984 & 1 & Atestação confirmada \cite{anual1984}. \\
1985 & 1 & Atestação confirmada \cite{camara1985}. \\
1986 & 1 & Atestação confirmada \cite{anual1986}. \\
1987 & 1 & Atestação confirmada \cite{anual1987}. \\
1988 & 1 & Atestação confirmada \cite{anual1988}. \\
1989 & 1 & Atestação confirmada \cite{anual1989}. \\
1990 & 1 & Atestação confirmada \cite{anual1990}. \\
1991 & 1 & Atestação confirmada \cite{anual1991}. \\
1992 & 1 & Atestação confirmada \cite{anual1992}. \\
1993 & 1 & Atestação confirmada \cite{anual1993}. \\
1994 & 1 & Atestação confirmada \cite{folha1994}. \\
1995 & 1 & Atestação confirmada \cite{anual1995}. \\
1996 & 1 & Atestação confirmada \cite{anual1996}. \\
1997 & 1 & Atestação confirmada \cite{anual1997}. \\
1998 & 1 & Atestação confirmada \cite{anual1998}. \\
1999 & 1 & Atestação confirmada \cite{anual1999}. \\
2000 & 1 & Atestação confirmada \cite{anual2000}. \\
2001 & 1 & Atestação confirmada \cite{anual2001}. \\
2002 & 1 & Atestação confirmada \cite{anual2002}. \\
2003 & 1 & Atestação confirmada \cite{anual2003}. \\
2004 & 1 & Atestação confirmada \cite{anual2004}. \\
2005 & 1 & Atestação confirmada \cite{anual2005}. \\
2006 & 1 & Atestação confirmada \cite{anual2006}. \\
2007 & 1 & Atestação confirmada \cite{young2008}. \\
2008 & 1 & Atestação confirmada \cite{anual2008}. \\
2009 & 1 & Atestação confirmada \cite{anual2009}. \\
\bottomrule\end{longtable}\normalsize

\section{Catálogo de pistas fornecidas para ampliação}
Este catálogo preserva as sugestões de jornais, revistas e obras que motivaram a investigação. É uma lista de pistas, não um conjunto de fontes que comprovam as grafias atribuídas. As descrições herdadas são alegações do roteiro; datas de composição, publicação, reedição e disponibilização digital precisam ser distinguidas. A presença de uma obra no catálogo não autentica uma passagem nem valida sua alocação decenal.
\small\begin{longtable}{p{0.09\textwidth}p{0.13\textwidth}p{0.66\textwidth}}
\caption{Pistas herdadas do roteiro, conservadas para auditoria. Descrições não constituem resultados verificados.}\\
\toprule Década alegada & Tipo & Descrição herdada (a verificar) \\\midrule\endfirsthead
\toprule Década alegada & Tipo & Descrição herdada (a verificar) \\\midrule\endhead
1900 & periódico & Gazeta de Notícias (RJ): Nas crônicas satíricas da virada do século, o riso irônico de salão aparecia rigidamente espaçado por pontuações: "Ah! Ah! Ah!". \\
1900 & periódico & O Paiz (RJ): Em artigos de crônica mundana relatando mexericos políticos do início da República. \\
1900 & periódico & Almanak Laemmert (RJ): Nas seções de anedotas e variedades comerciais, registrava o riso clássico pontuado. \\
1900 & obra & Canaã (1902) - Graça Aranha: O clássico do pré-modernismo utiliza a interjeição de forma solene e isolada ao longo dos diálogos. \\
1900 & obra & Recordações do Escrivão Isaías Caminha (1909) - Lima Barreto: Em diálogos que retratam a redação de um jornal carioca da época, os risos são marcados de forma elegante. \\
1900 & obra & A Esfinge (1908) - Coelho Neto: Romance que espelha os salões e teatros cariocas da época da Belle Époque, com diálogos tipograficamente aristocráticos. \\
1910 & periódico & Revista Careta (RJ): Pioneira no uso de piadas curtas ilustradas com a grafia de transição "Hê, hê, hê!". \\
1910 & periódico & Revista O Malho (RJ): O grande semanário satírico ilustrado registrava gargalhadas abertas em charges como "Quá... Quá...". \\
1910 & periódico & A Cigarra (SP): Crônicas paulistas misturando o riso satírico com a modernidade que nascia na capital do café. \\
1910 & obra & Triste Fim de Policarpo Quaresma (1911) - Lima Barreto: O deboche dos burocratas e militares contra o protagonista nacionalista é pontuado com o seco "Ah! ah!". \\
1910 & obra & Urupês (1918) - Monteiro Lobato: O autor introduz o linguajar caipira e regionalista no papel, trazendo a sonoridade crua de risos curtos em suas histórias. \\
1910 & obra & Eu (1912) - Augusto dos Anjos: Embora o livro seja de poesias poéticas e expressionistas, há o registro do riso escarninho da humanidade na literatura da época. \\
1920 & periódico & Correio da Manhã (RJ): Nas crônicas diárias registrando diálogos de salão e fofocas da elite carioca usando "Ah! ah! ah!". \\
1920 & periódico & Diário Nacional (SP): Órgão da imprensa modernista em São Paulo que usava o riso satírico de vanguarda. \\
1920 & periódico & Jornal do Recife (PE): Em seções e crônicas nordestinas. \\
1920 & obra & Macunaíma (1928) - Mário de Andrade: O herói sem nenhum caráter solta risadas folclóricas e descontraídas textualmente ao longo de suas aventuras. \\
1920 & obra & Memórias Sentimentais de João Miramar (1924) - Oswald de Andrade: O uso experimental da linguagem fragmentada reduz o riso a blocos rápidos. \\
1920 & obra & A Menina Morta (Escrito nos anos 20) - Cornélio Pena: Retrata a transição da fala interiorana para o papel com os risos teatrais pontuados. \\
1930 & periódico & Revista Fon-Fon! (RJ): O grande semanário cômico-cultural registrava o deboche de rua com o fonema "Quá! quá! quá!". \\
1930 & periódico & O Radical (RJ): Diário político combativo da época que ironizava adversários políticos fechando notas jornalísticas com escárnios onomatopaicos. \\
1930 & periódico & A Batalha (RJ): Outro veículo forte da imprensa satírica e combativa dos anos 30. \\
1930 & obra & O Quinze (1930) - Rachel de Queiroz: O romance de 30 traz diálogos crus do sertão, transcrevendo risos breves. \\
1930 & obra & São Bernardo (1934) - Graciliano Ramos: O riso cínico ou irônico dos fazendeiros é marcado de forma seca e direta na prosa. \\
1930 & obra & Reforma Agrária (1934) - Marquie: Ensaios e crônicas satíricas da política da década de 1930. \\
1940 & periódico & Diário de Notícias (RJ): Transcrições formais de entrevistas com a rubrica descritiva institucionalizada de (Risos) ou (Gargalhadas). \\
1940 & periódico & A Noite (RJ): Registro das transmissões de rádio ao vivo com rubricas técnicas indicando ao leitor o riso do auditório. \\
1940 & periódico & Folha da Manhã (SP): Registro sóbrio e formal do jornalismo paulistano dos anos 40. \\
1940 & obra & Fogo Morto (1943) - José Lins do Rego: O declínio dos engenhos de açúcar retratado com diálogos fortes e risos de escárnio social. \\
1940 & obra & O Tempo e o Vento: O Continente (1949) - Erico Verissimo: A saga épica gaúcha traz risos de caserna e de taberna descritos literariamente. \\
1940 & obra & Sagarana (1946) - João Guimarães Rosa: O experimentalismo mineiro usa onomatopeias independentes e poéticas para recriar o riso do sertanejo. \\
1950 & periódico & Revista do Pato Donald (1950 - Editora Abril): O berço que massificou o "Quá-quá-quá" com hífens no imaginário jovem brasileiro. \\
1950 & periódico & Tribuna da Imprensa (RJ): Liderado por Carlos Lacerda, o jornal usava o riso sarcástico nas colunas de denúncia política ("He-he-he..."). \\
1950 & periódico & Última Hora (RJ): Jornal moderno e popular de Samuel Wainer, registrando crônicas cotidianas e urbanas. \\
1950 & obra & Grande Sertão: Veredas (1956) - João Guimarães Rosa: Transcreve a fala e o riso dos jagunços através de sonoridades guturais únicas. \\
1950 & obra & Vestido de Noiva (Publicado em livro nos anos 50) - Nelson Rodrigues: O teatro rodriguiano introduz a gargalhada histérica ou maliciosa nas rubricas e falas suburbanas. \\
1950 & obra & Crônicas de Stanislaw Ponte Preta (Iniciadas nos anos 50) - Sérgio Porto: Crônica urbana carioca pré-bossa nova abusando do riso malicioso da noite de Copacabana. \\
1960 & periódico & O Pasquim (1969): O jornal contracultural que revolucionou a linguagem, agrupando letras fluidamente em suas famosas entrevistas descritas ("hahaha" ou "hehehe" soltos). \\
1960 & periódico & Jornal do Brasil (RJ): Com sua famosa reforma gráfica dos anos 60, modernizou a transcrição visual de entrevistas. \\
1960 & periódico & Revista Manchete (RJ): Cobertura da Jovem Guarda e do comportamento jovem usando onomatepeias corridas. \\
1960 & obra & Quarup (1967) - Antônio Callado: O romance da década de 60 traz diálogos descontraídos e informais da juventude e de intelectuais. \\
1960 & obra & A Paixão Segundo G.H. (1964) - Clarice Lispector: O fluxo de consciência foca no sorriso mental ou irônico introspectivo. \\
1960 & obra & Zero (Escrito nos anos 60) - Ignácio de Loyola Brandão: Romance experimental e fragmentado que usa a colagem de textos de rua e onomatopeias pop. \\
1970 & periódico & Gibi da Turma da Mônica (Anos 70 - Editora Abril): Consolidação do "Hi-hi-hi" hifenizado (riso escondido do Cebolinha) e do "Ha-ha-ha". \\
1970 & periódico & Opinião (SP): Jornal semanário de resistência política que usava ironias nas entrelinhas textuais. \\
1970 & periódico & Movimento (SP): Imprensa alternativa dos anos 70 capturando debates com marcações de diálogos. \\
1970 & obra & As Meninas (1973) - Lygia Fagundes Telles: O retrato de três jovens na ditadura usa o riso descontraído de campus universitário em seus diálogos. \\
1970 & obra & Sargento Getúlio (1971) - João Ubaldo Ribeiro: A risada seca e brutal do policial nordestino transcreve a sonoridade sertaneja. \\
1970 & obra & A Morte de D.J. em Paris (1975) - Roberto Drummond: Prosa pop e urbana influenciada diretamente pelas linguagens dos quadrinhos e da publicidade. \\
1980 & periódico & Folha de S.Paulo (SP): Nas colunas sociais cotidianas e bastidores da política da seção Painel, o termo "risos" colapsa e é registrado abreviado como "Rs." ou "Rsrs". \\
1980 & periódico & Jornal do Brasil (RJ): Uso comum do "rs" na coluna de fofocas e cultura informal de Zózimo Barrozo do Amaral. \\
1980 & periódico & O Globo (RJ): Notas rápidas de bastidores da televisão e celebridades adotando a economia mecânica da datilografia. \\
1980 & obra & Viva o Povo Brasileiro (1984) - João Ubaldo Ribeiro: O épico histórico reproduz o riso das variadas classes sociais brasileiras ao longo dos séculos. \\
1980 & obra & O Alquimista (1988) - Paulo Coelho: A literatura de massa dos anos 80 adota uma estrutura limpa e direta de diálogos com interjeições universais de riso ("rires", "sorrisos"). \\
1980 & obra & Feliz Ano Velho (1982) - Marcelo Rubens Paiva: O clássico juvenil dos anos 80 é escrito na exata linguagem coloquial da juventude paulistana da época, com risos corridos e espontâneos de conversa. \\
1990 & periódico & Revista Veja (Anos 90): Nas tradicionais Páginas Amarelas de entrevista, manteve o padrão formal de (Risos) entre parênteses para denotar a fala de políticos e artistas. \\
1990 & periódico & Folhateen (Suplemento jovem da Folha de S.Paulo): O berço impresso que documentava as gírias do mIRC e BBS, registrando graficamente o uso do "rsrsrs" e do "hahaha" que os jovens digitavam nas salas. \\
1990 & periódico & Revista Bizz (Mundo da Música): Jornalismo cultural informal capturando entrevistas com risos transcritos sem rigidez tipográfica. \\
1990 & obra & Chove nos Campos de Cachoeira (1998) - Dalcídio Jurandir: Estilos literários tardios que mantinham viva a tradição descritiva de diálogos no papel. \\
1990 & obra & Budapeste (Escrito/Planejado no fim dos anos 90) - Chico Buarque: Foca na minúcia das palavras e nos dilemas da tradução e grafia de expressões idiomáticas. \\
1990 & obra & Estorvo (1991) - Chico Buarque: Narrativa asfixiante e urbana em que os risos e o humor são tensos, espelhando a transição social dos anos 90. \\
2000 & periódico & Revista Capricho (Anos 2000): A publicação jovem foi uma das primeiras a colocar matérias impressas simulando telas de MSN, imprimindo pela primeira vez termos como "kakakaka" e "kkkkk" em papel. \\
2000 & periódico & Caderno de Informática (O Globo / Folha): Reportagens e colunas da época dedicadas a "traduzir o internetês" dos jovens, publicando tabelas que explicavam o que significavam o kkkk e o vc para os pais assustados. \\
2000 & periódico & Jornal Meia Hora (RJ - Fundado em 2005): Jornalismo popular de massa que quebrou a barreira da sobriedade e passou a adotar manchetes com a linguagem escrachada da internet dos anos 2000. \\
2000 & obra & Os Segredos do Orkut (2005) - Livros de nicho digital: Publicações focadas em desvendar as primeiras grandes redes sociais brasileiras, documentando os logs de depoimentos repletos de "kkkkk". \\
2000 & obra & Manual do MSN Messenger (Meados dos anos 2000): Guias práticos em livrarias ensinando a configurar a ferramenta e a decodificar as abreviações e os novos emoticons rindo. \\
2000 & obra & Literatura de Blog (Coletâneas dos anos 2000): O nascimento de livros impressos baseados em posts de blogs famosos da época (como os textos de Cardosonline ou Melhores do Mundo), que transportaram o "kkkkk" nativo das telas diretamente para o catálogo das editoras brasileiras. \\
\bottomrule\end{longtable}\normalsize

\section{Parecer individual da segunda auditoria}
Os identificadores R01--R66 seguem a mesma ordem do catálogo anterior. A categoria A exige representação do riso autenticada; L identifica referência lexical, temática ou cotejo contextual, sem autenticar a grafia atribuída; B registra identificação bibliográfica; I registra recuperação do exemplar específico sem autenticação da frase; S conserva casos sem ocorrência autenticada ou com identificação insuficiente. Há três A, sete L, dezenove B, quatro I e trinta e três S. Esses estados descrevem o trabalho documental, não a frequência histórica de qualquer grafia. URLs, estado de aquisição e data do parecer constam no CSV público.
\small\begin{longtable}{p{0.07\textwidth}p{0.08\textwidth}p{0.72\textwidth}}
\caption{Parecer individual sobre as 66 indicações. A: representação autenticada; L: referência lexical; B: identificação bibliográfica; I: exemplar específico; S: sem ocorrência autenticada.}\\
\toprule ID & Estado & Resultado e limites \\\midrule\endfirsthead
\toprule ID & Estado & Resultado e limites \\\midrule\endhead
R01 & I & Exemplar de 25/07/1904, n. 207, recuperado; busca no OCR das quatro páginas não autenticou a frase do ovo. OCR degradado impede conclusão negativa definitiva.\\
R02 & S & Busca por título, década e riso não individualizou artigo, edição e página. A descrição genérica não comprova grafia.\\
R03 & S & Sem ano e edição do almanaque ou título da anedota. Busca não individualizou ocorrência; é necessário delimitar volume e seção.\\
R04 & L & Texto integral recuperado: expressão de riso, PDF p. 28. Referência lexical; não autentica o diálogo atribuído ao colono Maciel.\\
R05 & L & Exemplar consultado: segunda edição, 1917, PDF p. 26: Riu-se muito. Não datar a passagem de 1909 sem cotejo da primeira edição.\\
R06 & L & Oito capítulos consultados; capítulo VI contém risos casquinados. É descrição lexical, não uma sequência onomatopaica demonstrada.\\
R07 & S & A busca recuperou Quá em Quase doutor, Careta, 08/05/1915, via Vida urbana. Contexto de recusa, não riso: falso positivo excluído.\\
R08 & I & Edição n. 526 de 12/10/1912 confirmada na capa. Busca nas 58 páginas do OCR não autenticou prefeito/carne barata; desenhos exigem revisão visual integral.\\
R09 & S & Pista alternativa: poema Epístola, A Cigarra, n. 16, 20/01/1915, p. 27, em reprodução externa. Não conferida no exemplar original; não incluída no corpus.\\
R10 & L & Texto de 1915 consultado. Terceira parte, cap. II, contém Você, Quaresma, é um visionario; sem Ah! ah! e sem a continuação sobre plantar. Citação recebida não é transcrição fiel.\\
R11 & A & Quinta edição de 1919: p. impressa 19, PDF p. 24: Quá, quá, quá!, com risos explícitos no contexto. Não retrodatado automaticamente a 1918.\\
R12 & L & Monólogo de uma sombra: Chorando e rindo na ironia infausta. Confirma referência temática, não onomatopeia.\\
R13 & S & Frase de 14/02/1925 não autenticada. Fon-Fon!, maio de 1925, p. 3, noticia retorno do Correio após suspensão; é indício de problema cronológico, não prova autônoma de que não existiu edição em fevereiro.\\
R14 & A & 17/06/1928, n. 290, p. 4: (risos, muito bem), em discurso de Mattos Pimenta. Rubrica descritiva, não grafia de gargalhada.\\
R15 & S & Busca por título, década e risos não individualizou reportagem verificável; faltam dia e página.\\
R16 & A & Edição UFFS de 2019, p. 139: Quá! quá! quá quaquá!... e indicação de gargalhada. Obra original de 1928; a passagem não é a frase Ah! ah! enviada.\\
R17 & L & Exemplar de 1924, p. 114, PDF p. 116: morreu nos lábios de todos o sorriso. Referência lexical; não comprova onomatopeia.\\
R18 & B & BBM-USP registra edição de 1954. A alegação de composição nos anos 1920 não foi demonstrada. Nenhuma passagem de riso autenticada.\\
R19 & I & PDF identificado pelo acervo como n. 23 de 1935 recuperado; capa e OCR consultados. Data de 08/06 e frase sobre rei da Inglaterra não autenticadas; falta conferir expediente legível.\\
R20 & S & Exemplar n. 1175 de 1935 recuperado e OCR pesquisado. Há candidatos lexicais, mas nenhuma grafia atribuída autenticada por imagem.\\
R21 & S & Exemplar n. 31 de 1930 recuperado e OCR pesquisado. Referências lexicais candidatas não autenticam a generalização sobre marcações de riso.\\
R22 & B & Editora Record confirma estreia em 1930; aquisição automática recebeu HTTP 429. Não se confirmou passagem nem grafia.\\
R23 & L & Provas da segunda edição, 1938, recuperadas; PDF p. 63: sorrisos frios. A frase latim de missa não foi localizada no OCR; não retrodatado a 1934.\\
R24 & S & Reforma Agrária (1934), de Marquie: buscas por título, ano e variante do nome não identificaram registro bibliográfico suficiente. Não tratar como obra inexistente.\\
R25 & I & Edição de 18/03/1944, n. 6564, confirmada no cabeçalho. Busca do OCR não autenticou ministro/boatos de racionamento; texto de baixa qualidade.\\
R26 & S & Busca por título, ano e rubrica não individualizou transmissão ou reportagem com página. Não demonstrada a descrição recebida.\\
R27 & S & Busca produziu falso positivo: Jornal de Notícias, de 02/03/1947, e não Folha da Manhã. Documento incluído separadamente; não substitui esta referência.\\
R28 & B & Bibliografia da Academia Brasileira de Letras confirma Fogo morto, 1943. Texto e passagem sobre riso não conferidos.\\
R29 & B & Editora identifica O Continente na trilogia publicada entre 1949 e 1962. Metadado não autentica diálogo de caserna ou taberna.\\
R30 & B & Museu Casa Guimarães Rosa registra exemplar de Sagarana de 1946. Frase He! he! e mundo dá muitas voltas não autenticada.\\
R31 & S & Revista de 1950 e alegado Almanaque de 1953 não são a mesma referência. Faltam número, história e página para verificar moedas e Quá-quá-quá.\\
R32 & S & Busca de frase e jornal não localizou plano é honesto. Ano de 1955 sem dia, coluna e página não individualiza o documento.\\
R33 & S & Busca por jornal, década e risos não individualizou crônica. É necessária referência de edição.\\
R34 & B & Página editorial de Grande sertão: veredas consultada; primeira edição em 1956. Nenhuma passagem específica de riso foi autenticada.\\
R35 & B & Registro institucional sobre Vestido de noiva identifica escrita em 1943; estreia e publicação são eventos distintos. A edição em livro dos anos 1950 não foi individualizada.\\
R36 & S & Crônicas de Stanislaw Ponte Preta não designa título, edição ou crônica específica. Identidade do cronista não confirma data e grafia de passagem.\\
R37 & S & Coleção identificada no dossiê da Biblioteca Nacional. URL tentada para n. 12 retornou 404; diálogo atribuído a Jaguar permanece não autenticado.\\
R38 & S & A reforma gráfica do jornal não demonstra uso de hahaha. Busca não localizou entrevista e página específicas na década de 1960.\\
R39 & S & Sem edição e entrevistado identificados. Busca por título, década e rubrica não autenticou onomatopeia corrida.\\
R40 & B & Página editorial de Quarup localizada; download automático HTTP 429. Frase sobre jornais não autenticada e sem página.\\
R41 & B & Registro museológico localizado para A paixão segundo G.H.; download retorna página sem conteúdo textual suficiente. Não confirma sorriso mental ou grafia.\\
R42 & B & ABL distingue edição italiana de 1974 e portuguesa de 1975. Alegada composição nos anos 1960 não autoriza contar publicação naquela década.\\
R43 & S & Menção genérica a gibis dos anos 1970. Busca não individualizou número e história com as formas atribuídas.\\
R44 & S & Busca por Opinião, década e risos não individualizou passagem verificável. A caracterização ideológica não prova grafia.\\
R45 & S & Busca por Movimento, década e risos não individualizou debate com localizador. Pendente identificação de edição.\\
R46 & B & Página editorial de As meninas consultada. Frase sobre aposta e Ha-ha-ha! não autenticada; sem acesso integral conferido.\\
R47 & B & Página editorial localizada na busca; aquisição automática falhou com 404. Data de obra não é prova de uma risada específica.\\
R48 & B & Coletânea de 1975 identificada também em pesquisa institucional UFRN; catálogo WorldCat bloqueou aquisição com 403. Nenhuma passagem de riso confirmada.\\
R49 & S & Busca da data 14/11/1986 e frase de congelamento não autenticou Rs. na seção Painel. Rubrica risos de outra publicação de 1994 não substitui esta alegação.\\
R50 & S & Busca por Zózimo, década e rs não individualizou nota. Não demonstrado uso comum de rs no jornal nos anos 1980.\\
R51 & S & Busca por jornal, década e rsrs não individualizou nota. Alegada economia da datilografia permanece explicação sem evidência específica.\\
R52 & B & Página editorial de Viva o povo brasileiro consultada. Descrição literária genérica não comprova forma particular de riso.\\
R53 & B & Página editorial de O Alquimista consultada. Sorriso e rir são itens lexicais; não devem ser classificados automaticamente como interjeições.\\
R54 & B & Página editorial de Feliz ano velho consultada. Frase burradas do colégio e sequência Ha ha ha não autenticadas.\\
R55 & S & Busca de Veja, 07/05/1993 e frase não fugirei da raia não individualizou entrevista. Data e entrevistado continuam sem autenticação.\\
R56 & S & Busca de Folhateen, junho de 1997, mIRC e grafias não identificou reportagem com dia e página. Não confirmada a frase recebida.\\
R57 & S & Busca de Bizz, década e hahaha não identificou entrevista específica. Sem edição, autoria e página.\\
R58 & B & Fonte institucional UEPA registra primeira publicação em 1941. Ano de 1998 pode designar edição; não é data de nascimento da obra nem da passagem.\\
R59 & B & Página editorial informa lançamento em 10/09/2003. Planejamento nos anos 1990 não demonstrado e não equivale a publicação.\\
R60 & B & Página editorial informa lançamento em 02/08/1991. Nenhuma passagem sobre riso foi conferida no texto integral.\\
R61 & S & Busca de Capricho, setembro de 2007 e frase quase quebrou a cara não individualizou reportagem. Não se sustenta afirmação de primeira impressão de kkkk.\\
R62 & S & Caderno de Informática de dois jornais é referência ambígua. Busca por internetês e kkkk não individualizou tabela de abreviações com data e página.\\
R63 & S & Busca por Meia Hora, 2005 e kkkk não identificou manchete. Fundação do veículo não demonstra adoção de grafia específica.\\
R64 & S & Título, autor e editora do alegado livro de 2005 não identificados. Resultado homônimo é postagem de blog de 2008, que não autentica o livro.\\
R65 & S & Manual do MSN Messenger sem autor, ISBN ou editora. Busca não individualizou obra e edição; não autentica convenções descritas.\\
R66 & S & Literatura de Blog é categoria, não referência bibliográfica individual. Cardosonline e Melhores do Mundo não podem ser convertidos em livro sem título, autor e edição.\\
\bottomrule\end{longtable}\normalsize

\section{Correções conceituais e decisões editoriais}
\begin{enumerate}[leftmargin=*]
\item Mudança histórica e motivação icônica são dimensões distintas, que podem coexistir.
\item Continuidade lexical não exige afixação produtiva ou radical superficialmente preservado.
\item Uma cadeia histórica pode ser reconstruída sem documentação de todas as gerações. A ausência de elos observados reduz a força da demonstração, mas não prova independência.
\item Não há prova de que \forma{quiá} desapareceu da circulação social ao ser impresso. A extensão de uso requer corpus.
\item \forma{Quiá} em \textit{Til} não é ocorrência literal de \forma{kkkk}, nem demonstra a primeira representação de riso da língua.
\item \forma{Você} para \forma{cê} envolve aférese. A cronologia que situa \forma{você} apenas no século XX contradiz a atestação de 1638 discutida por Monte.
\item \forma{Cá} é nome de letra ou sílaba, não um fonema único; /b/ é bilabial; a gargalhada não contém necessariamente /k/.
\item \forma{Vilão} não se explica como simples mudança de \forma{-us} em \forma{-ão}. Mudança semântica não determina sozinha classe gramatical.
\item Famílias linguísticas diferentes não excluem empréstimo; \forma{chikutaku} não é prova de isolamento.
\item Datas de plataformas, reformas ortográficas e teclados são contexto potencial, não provas de origem de uma sequência.
\item Preferências geracionais não são universais; o corpus não estima usos brasileiros da geração Alpha.
\item Uma referência de divulgação citada num artigo acadêmico não se torna novo testemunho primário.
\item A quota de três jornais e três livros é editorial. A justificativa estatística depende do estimando e da coleta.
\item Datas precisas e aspas exigem localizador e conferência. Ilustrações não podem figurar como transcrições.
\end{enumerate}

\section{Dicionário de dados e reprodução}
\texttt{attestations.csv} reúne testemunhos; \texttt{claims\_audit.csv}, os 22 alvos; \texttt{annual\_coverage.csv}, a cobertura anual; \texttt{candidate\_sources.csv}, as pistas; \texttt{downloads.json}, as aquisições; \texttt{detection\_sensitivity.csv}, os cenários matemáticos; \texttt{summary.json}, as contagens.

\texttt{build\_data.py} gera tabelas de registros codificados e parâmetros declarados. A aquisição é separada da análise para que falhas de rede não sejam interpretadas como resultados linguísticos. Regex localizam candidatos e não substituem anotação. A segunda rodada acrescenta os arquivos \texttt{references\_audit.csv}, \texttt{claims\_verification.csv} e \texttt{extra\_attestations.json}, além dos registros de buscas. O repositório documenta comandos e versões.

Não foram produzidos logs sintéticos para os anos 1990, participantes fictícios, contagens populacionais ou intervalos de confiança para a amostra dirigida. Os cenários de planejamento são identificados como hipotéticos. A distinção entre realizado e proposto integra a metodologia.

\end{document}
